\documentclass[oneside]{book}

% ------------------------------------------------------------------------------
% Packages
% ------------------------------------------------------------------------------

\usepackage[utf8]{inputenc}
\usepackage[english]{babel}

\usepackage{amsmath, amsthm, amssymb, amsfonts}
\usepackage{mathtools}

\usepackage[dvipsnames]{xcolor}
\usepackage{tcolorbox}
\tcbuselibrary{theorems, skins, breakable}


\usepackage{geometry}
\usepackage{setspace}
\usepackage{hyperref}
\usepackage{microtype}
\usepackage{booktabs}
\usepackage{array}

\usepackage{multicol}
\setlength{\columnsep}{1.2em}

\usepackage{graphicx}

% ------------------------------------------------------------------------------
% Page Setup
% ------------------------------------------------------------------------------

\geometry{
    top=0.8in,
    bottom=0.8in,
    left=0.7in,
    right=0.7in,
    headheight=12pt,
    headsep=20pt,
    footskip=25pt
}

\setstretch{1.15}

\hypersetup{
    colorlinks=true,
    linkcolor=MidnightBlue,
    urlcolor=RoyalBlue,
    citecolor=ForestGreen
}

% ------------------------------------------------------------------------------
% Variables
% ------------------------------------------------------------------------------

\def\notetitle{Machine Learning: A Field Guide for the Perpetually Confused}
\def\notesubtitle{From ``What Is a Gradient'' to ``Why Is My Loss Exploding''}
\def\noteauthor{Erdan Mu Alan}
\def\noteinstitution{University of Pennsylvania}

% ------------------------------------------------------------------------------
% Theorem System and User Commands
% ------------------------------------------------------------------------------

% !TEX root = main.tex

% Theorem System
% The following boxes are provided:
%   Definition:     \defn 
%   Theorem:        \thm 
%   Lemma:          \lem
%   Corollary:      \cor
%   Proposition:    \prop   
%   Claim:          \clm
%   Fact:           \fact
%   Proof:          \pf
%   Example:        \ex
%   Remark:         \rmk (sentence), \rmkb (block)
% Suffix
%   r:              Allow Theorem/Definition to be referenced, e.g. thmr
%   p:              Add a short proof block for Lemma, Corollary, Proposition or Claim, e.g. lemp
%                   For theorems, use \pf for proof blocks

\tcbset{
    kb/common block/.style={
        enhanced,
        breakable,
        frame hidden,
        boxrule=0pt,
        sharp corners,
        boxsep=0pt,
        left=6mm,
        right=5mm,
        top=4mm,
        bottom=4mm,
        before skip=10pt,
        after skip=12pt,
        before upper={\parindent15pt\relax},
    },
    kb/proof block/.style={
        kb/common block,
        colback=black!1!white,
        borderline west={0.9mm}{0pt}{black!55},
    },
    kb/accent proof/.style n args={2}{
        kb/common block,
        colback=#1!4!white,
        borderline west={1mm}{0pt}{#1},
        before upper={
            \parindent15pt\relax
            {\noindent\bfseries #2}\par\smallskip
        },
    },
    kb/example block/.style={
        kb/common block,
        colback=MidnightBlue!3!white,
        borderline west={1.2mm}{0pt}{MidnightBlue!70!black},
        borderline north={0.3mm}{0pt}{MidnightBlue!18!white},
    },
    kb/example title/.style={
        colback=MidnightBlue!75!black,
        colframe=MidnightBlue!75!black,
        boxrule=0pt,
        sharp corners,
        left=3mm,
        right=3mm,
        top=1.2mm,
        bottom=1.2mm,
    },
}

\newcounter{example}[section]
\renewcommand{\theexample}{\thesection.\arabic{example}}

% Definition
\newtcbtheorem[number within=section]{mydefinition}{Definition}
{
    enhanced,
    frame hidden,
    titlerule=0mm,
    toptitle=1mm,
    bottomtitle=1mm,
    fonttitle=\bfseries\large,
    coltitle=black,
    colbacktitle=blue!20!white,
    colback=blue!10!white,
}{defn}

\NewDocumentCommand{\defn}{m+m}{
    \begin{mydefinition}{#1}{}
        #2
    \end{mydefinition}
}

\NewDocumentCommand{\defnr}{mm+m}{
    \begin{mydefinition}{#1}{#2}
        #3
    \end{mydefinition}
}

\NewDocumentEnvironment{definition}{o+b}{
    \IfNoValueTF{#1}{
        \begin{mydefinition}{}{}
    }{
        \begin{mydefinition}{#1}{}
    }
    #2
    \end{mydefinition}
}{}

% Theorem
\newtcbtheorem[use counter from=mydefinition]{mytheorem}{Theorem}
{
    enhanced,
    frame hidden,
    titlerule=0mm,
    toptitle=1mm,
    bottomtitle=1mm,
    fonttitle=\bfseries\large,
    coltitle=black,
    colbacktitle=cyan!20!white,
    colback=cyan!10!white,
}{thm}

\NewDocumentCommand{\thm}{m+m}{
    \begin{mytheorem}{#1}{}
        #2
    \end{mytheorem}
}

\NewDocumentCommand{\thmr}{mm+m}{
    \begin{mytheorem}{#1}{#2}
        #3
    \end{mytheorem}
}

\NewDocumentEnvironment{theorem}{o+b}{
    \IfNoValueTF{#1}{
        \begin{mytheorem}{}{}
    }{
        \begin{mytheorem}{#1}{}
    }
    #2
    \end{mytheorem}
}{}

% Lemma
\newtcbtheorem[use counter from=mydefinition]{mylemma}{Lemma}
{
    enhanced,
    frame hidden,
    titlerule=0mm,
    toptitle=1mm,
    bottomtitle=1mm,
    fonttitle=\bfseries\large,
    coltitle=black,
    colbacktitle=violet!20!white,
    colback=violet!10!white,
}{lem}

\NewDocumentCommand{\lem}{m+m}{
    \begin{mylemma}{#1}{}
        #2
    \end{mylemma}
}

\newenvironment{lempf}{
    \begin{tcolorbox}[kb/accent proof={violet!65!black}{Proof for Lemma.}]
}{
    \hfill\textcolor{violet!65!black}{$\blacksquare$}
    \end{tcolorbox}
}

\NewDocumentCommand{\lemp}{m+m+m}{
    \begin{mylemma}{#1}{}
        #2
    \end{mylemma}

    \begin{lempf}
        #3
    \end{lempf}
}

% Corollary
\newtcbtheorem[use counter from=mydefinition]{mycorollary}{Corollary}
{
    enhanced,
    frame hidden,
    titlerule=0mm,
    toptitle=1mm,
    bottomtitle=1mm,
    fonttitle=\bfseries\large,
    coltitle=black,
    colbacktitle=orange!20!white,
    colback=orange!10!white,
}{cor}

\NewDocumentCommand{\cor}{+m}{
    \begin{mycorollary}{}{}
        #1
    \end{mycorollary}
}

\newenvironment{corpf}{
    \begin{tcolorbox}[kb/accent proof={orange!80!black}{Proof for Corollary.}]
}{
    \hfill\textcolor{orange!80!black}{$\blacksquare$}
    \end{tcolorbox}
}

\NewDocumentCommand{\corp}{m+m+m}{
    \begin{mycorollary}{}{}
        #1
    \end{mycorollary}

    \begin{corpf}
        #2
    \end{corpf}
}

% Proposition
\newtcbtheorem[use counter from=mydefinition]{myproposition}{Proposition}
{
    enhanced,
    frame hidden,
    titlerule=0mm,
    toptitle=1mm,
    bottomtitle=1mm,
    fonttitle=\bfseries\large,
    coltitle=black,
    colbacktitle=yellow!30!white,
    colback=yellow!20!white,
}{prop}

\NewDocumentCommand{\prop}{+m}{
    \begin{myproposition}{}{}
        #1
    \end{myproposition}
}

\newenvironment{proppf}{
    \begin{tcolorbox}[kb/accent proof={Goldenrod!85!black}{Proof for Proposition.}]
}{
    \hfill\textcolor{Goldenrod!85!black}{$\blacksquare$}
    \end{tcolorbox}
}

\NewDocumentCommand{\propp}{+m+m}{
    \begin{myproposition}{}{}
        #1
    \end{myproposition}

    \begin{proppf}
        #2
    \end{proppf}
}

% Claim
\newtcbtheorem[use counter from=mydefinition]{myclaim}{Claim}
{
    enhanced,
    frame hidden,
    titlerule=0mm,
    toptitle=1mm,
    bottomtitle=1mm,
    fonttitle=\bfseries\large,
    coltitle=black,
    colbacktitle=pink!30!white,
    colback=pink!20!white,
}{clm}


\NewDocumentCommand{\clm}{m+m}{
    \begin{myclaim*}{#1}{}
        #2
    \end{myclaim*}
}

\newenvironment{clmpf}{
    \begin{tcolorbox}[kb/accent proof={Magenta!70!black}{Proof for Claim.}]
}{
    \hfill\textcolor{Magenta!70!black}{$\blacksquare$}
    \end{tcolorbox}
}

\NewDocumentCommand{\clmp}{m+m+m}{
    \begin{myclaim*}{#1}{}
        #2
    \end{myclaim*}

    \begin{clmpf}
        #3
    \end{clmpf}
}

% Fact
\newtcbtheorem[use counter from=mydefinition]{myfact}{Fact}
{
    enhanced,
    frame hidden,
    titlerule=0mm,
    toptitle=1mm,
    bottomtitle=1mm,
    fonttitle=\bfseries\large,
    coltitle=black,
    colbacktitle=purple!20!white,
    colback=purple!10!white,
}{fact}

\NewDocumentCommand{\fact}{+m}{
    \begin{myfact}{}{}
        #1
    \end{myfact}
}


% Proof
\RenewDocumentEnvironment{proof}{O{\proofname}}{
    \par
    \pushQED{\qed}
    \begin{tcolorbox}[kb/proof block]
        {\noindent\bfseries #1}\par\smallskip
        \ignorespaces
}{
    \popQED
    \end{tcolorbox}
    \par
}

\let\proofenvironment\proof
\let\endproofenvironment\endproof

\makeatletter
\renewcommand{\proof}{
    \@ifnextchar\bgroup{\proofcommand}{\proofenvironment}
}
\makeatother

\NewDocumentCommand{\proofcommand}{+m}{
    \proofenvironment
        #1
    \endproofenvironment
}

\NewDocumentCommand{\pf}{o+m}{
    \IfNoValueTF{#1}{
        \begin{proof}
            #2
        \end{proof}
    }{
        \begin{proof}[#1]
            #2
        \end{proof}
    }
}

% Example
\NewDocumentEnvironment{example}{o+b}{
    \refstepcounter{example}
    \IfNoValueTF{#1}{
        \def\exampletitle{Example \theexample}
    }{
        \def\exampletitle{Example \theexample \textbar\ #1}
    }
    \begin{tcolorbox}[
        kb/example block,
        title={\exampletitle},
        fonttitle=\bfseries\small,
        coltitle=white,
        boxed title style={kb/example title},
        attach boxed title to top left={xshift=4mm,yshift*=-\tcboxedtitleheight/2},
        top=7mm,
    ]
        #2
    \end{tcolorbox}
}{}

\NewDocumentCommand{\ex}{+m}{
    \begin{example}
        #1
    \end{example}
}


% Remark
\NewDocumentCommand{\rmk}{+m}{
    {\it \color{blue!50!white}#1}
}

\newenvironment{remark}{
    \par
    \vspace{5pt}
    \begin{minipage}{\textwidth}
        {\par\noindent{\textbf{Remark.}}}
        \tcolorbox[blanker,breakable,left=5mm,
        before skip=10pt,after skip=10pt,
        borderline west={1mm}{0pt}{cyan!10!white}]
}{
        \endtcolorbox
    \end{minipage}
    \vspace{5pt}
}

\NewDocumentCommand{\rmkb}{+m}{
    \begin{remark}
        #1
    \end{remark}
}

% !TEX root = main.tex

\newcommand{\R}{\mathbb{R}}
\newcommand{\N}{\mathbb{N}}
\newcommand{\vect}[1]{\mathbf{#1}}
\newcommand{\norm}[1]{\left\lVert #1 \right\rVert}
\newcommand{\dd}{\mathop{}\!\mathrm{d}}
\newcommand{\grad}{\nabla}
\newcommand{\diver}{\operatorname{div}}
\newcommand{\pd}[2]{\frac{\partial #1}{\partial #2}}
\newcommand{\lcm}{\operatorname{lcm}}


% ------------------------------------------------------------------------------
% Document
% ------------------------------------------------------------------------------

\begin{document}

% ------------------------------------------------------------------------------
% Title Page
% ------------------------------------------------------------------------------

\begin{titlepage}
    \centering

    \vspace*{2.2cm}

    {\Huge\bfseries \notetitle \par}

    \vspace{0.9cm}

    \resizebox{0.88\textwidth}{!}{%
        \itshape \notesubtitle
    }

    \vspace{2.4cm}

    \vspace{1.2cm}

    {\large \noteauthor \par}

    \vspace{0.25cm}

    {\large \noteinstitution \par}

    \vfill
\end{titlepage}

% ------------------------------------------------------------------------------
% Front Matter
% ------------------------------------------------------------------------------

\frontmatter
\setcounter{tocdepth}{1}
\tableofcontents
\newpage

% ------------------------------------------------------------------------------
% Main Matter
% ------------------------------------------------------------------------------

\mainmatter

% ==============================================================================
% PART I: FOUNDATIONS
% ==============================================================================

\part{Foundations}

% ------------------------------------------------------------------------------
\chapter{The Learning Problem}
\label{ch:learnproblem}
% ------------------------------------------------------------------------------

\section{What Is Machine Learning?}

\defn{Machine Learning}{
A computer program is said to \textbf{learn} from experience $E$ with respect to some class of tasks $T$ and performance measure $P$, if its performance at tasks in $T$, as measured by $P$, improves with experience $E$. \rmk{(Mitchell, 1997)}
}

Every ML system sits at the intersection of three components:
\begin{itemize}
    \item \textbf{Model} --- the family of functions $f_\theta : \mathcal{X} \to \mathcal{Y}$ parameterized by $\theta$.
    \item \textbf{Data} --- observations $\mathcal{D} = \{(\vect{x}_i, y_i)\}_{i=1}^N$ used to constrain $\theta$.
    \item \textbf{Optimization} --- the algorithm that searches for the best $\theta$ given $\mathcal{D}$.
\end{itemize}

The three learning paradigms differ in the structure of data:
\begin{description}
    \item[Supervised learning] Labeled pairs $(\vect{x}, y)$. Goal: predict $y$ from $\vect{x}$.
    \item[Unsupervised learning] Unlabeled observations $\{\vect{x}_i\}$. Goal: discover structure.
    \item[Reinforcement learning] Agent--environment interaction with reward signals. Goal: find a policy maximizing cumulative reward.
\end{description}

\section{The Statistical Learning Setup}

\defn{Hypothesis Class}{
A \textbf{hypothesis class} $\mathcal{H}$ is a set of candidate functions $h: \mathcal{X} \to \mathcal{Y}$. Learning means searching for the best $h \in \mathcal{H}$.
}

\defn{Loss Function}{
A \textbf{loss function} $\ell: \mathcal{Y} \times \mathcal{Y} \to \R_{\geq 0}$ quantifies the cost of predicting $\hat{y}$ when the truth is $y$. Common choices:
\begin{center}
\begin{tabular}{lll}
\toprule
Name & Formula & Use case \\
\midrule
Squared loss & $\ell(\hat{y}, y) = (\hat{y} - y)^2$ & Regression \\
0-1 loss & $\ell(\hat{y}, y) = \mathbf{1}[\hat{y} \neq y]$ & Classification \\
Cross-entropy & $\ell(\hat{p}, y) = -y\log\hat{p} - (1-y)\log(1-\hat{p})$ & Binary classification \\
\bottomrule
\end{tabular}
\end{center}
}

\defn{Risk and Empirical Risk Minimization}{
The \textbf{true risk} (generalization error) of $h$ is the expected loss under the data distribution:
\[
R(h) = \mathbb{E}_{(\vect{x},y) \sim \mathcal{D}}[\ell(h(\vect{x}),\, y)].
\]
Since $\mathcal{D}$ is unknown, we minimize the \textbf{empirical risk} on the training set:
\[
\hat{R}(h) = \frac{1}{N}\sum_{i=1}^{N}\ell(h(\vect{x}_i),\, y_i).
\]
\textbf{Empirical Risk Minimization (ERM):} $\hat{h} = \arg\min_{h \in \mathcal{H}}\hat{R}(h)$.
}

The difference $R(\hat{h}) - \hat{R}(\hat{h})$ is the \textbf{generalization gap}. A model \textbf{overfits} if this gap is large; it \textbf{underfits} if both $R$ and $\hat{R}$ are large. The capacity of $\mathcal{H}$ controls this tradeoff.

\rmkb{
Data is divided into three non-overlapping sets:
\begin{description}
    \item[Training set] Used to minimize $\hat{R}(h)$.
    \item[Validation set] Used for model selection and hyperparameter tuning.
    \item[Test set] Used once, at the very end, as an unbiased estimate of $R(\hat{h})$.
\end{description}
Never touch the test set until final evaluation.
}

% ------------------------------------------------------------------------------
\chapter{Probability and Statistics}
% ------------------------------------------------------------------------------

\section{Random Variables}

\defn{Random Variable}{
A \textbf{random variable} $X$ is a function mapping outcomes from a probability space to $\R$. We write $P(X = x)$ or $P(X \leq x)$ for event probabilities.

For a \textbf{discrete} RV: the \textbf{PMF} $p_X(x) = P(X = x)$ satisfies $\sum_x p_X(x) = 1$; the \textbf{CDF} $F_X(x) = \sum_{x' \leq x} p_X(x')$.

For a \textbf{continuous} RV: the \textbf{PDF} $f_X(x) = \frac{dF_X(x)}{dx}$ satisfies $\int f_X(x)\,dx = 1$; note $f_X(x) \neq P(X = x)$. The \textbf{CDF} $F_X(x) = \int_{-\infty}^x f_X(t)\,dt$.
}

\section{Expectation and Variance}

\defn{Expectation}{
\[
E[X] = \begin{cases} \displaystyle\sum_x x\,p_X(x) & \text{(discrete)} \\[6pt] \displaystyle\int_{-\infty}^{\infty} x\,f_X(x)\,dx & \text{(continuous)} \end{cases}
\]
More generally, $E[g(X)] = \int g(x)\,f_X(x)\,dx$.
}

\prop{
\textbf{Linearity of expectation}: for any constants $a, b$ and random variables $X, Y$:
\[
E[aX + b] = aE[X] + b, \qquad E[X + Y] = E[X] + E[Y].
\]
Linearity holds regardless of whether $X$ and $Y$ are independent.
}

\defn{Variance}{
\[
\mathrm{Var}[X] = E\!\left[(X - E[X])^2\right] = E[X^2] - (E[X])^2.
\]
The \textbf{standard deviation} is $\sigma_X = \sqrt{\mathrm{Var}[X]}$.
}

\prop{
Key variance identities: $\mathrm{Var}[a] = 0$, $\mathrm{Var}[aX] = a^2\,\mathrm{Var}[X]$, and
\[
\mathrm{Var}[X + Y] = \mathrm{Var}[X] + \mathrm{Var}[Y] + 2\,\mathrm{Cov}[X,Y].
\]
If $X, Y$ are independent, $\mathrm{Var}[X + Y] = \mathrm{Var}[X] + \mathrm{Var}[Y]$.
}

\section{Common Distributions}

\begin{center}
\renewcommand{\arraystretch}{1.4}
\begin{tabular}{lp{5.8cm}cc}
\toprule
Distribution & PMF / PDF & Mean & Variance \\
\midrule
$\mathrm{Bernoulli}(p)$ & $p^x(1-p)^{1-x}$, $x\in\{0,1\}$ & $p$ & $p(1-p)$ \\
$\mathrm{Binomial}(n,p)$ & $\binom{n}{x}p^x(1-p)^{n-x}$, $x\in\{0,\ldots,n\}$ & $np$ & $np(1-p)$ \\
$\mathrm{Geometric}(p)$ & $p(1-p)^{k-1}$, $k=1,2,\ldots$ & $\tfrac{1}{p}$ & $\tfrac{1-p}{p^2}$ \\
$\mathrm{Poisson}(\lambda)$ & $\tfrac{e^{-\lambda}\lambda^k}{k!}$, $k=0,1,\ldots$ & $\lambda$ & $\lambda$ \\
$\mathrm{Uniform}(a,b)$ & $\tfrac{1}{b-a}$ for $x\in(a,b)$ & $\tfrac{a+b}{2}$ & $\tfrac{(b-a)^2}{12}$ \\
$\mathcal{N}(\mu,\sigma^2)$ & $\tfrac{1}{\sigma\sqrt{2\pi}}\exp\!\left(-\tfrac{(x-\mu)^2}{2\sigma^2}\right)$ & $\mu$ & $\sigma^2$ \\
$\mathrm{Exponential}(\lambda)$ & $\lambda e^{-\lambda x}$, $x \geq 0$ & $\tfrac{1}{\lambda}$ & $\tfrac{1}{\lambda^2}$ \\
\bottomrule
\end{tabular}
\end{center}

\section{Maximum Likelihood Estimation}

Given data $\{x_i\}_{i=1}^N$ assumed i.i.d.\ from $p(x;\,\theta)$, MLE finds the parameter making observations most probable.

\defn{Likelihood, Log-Likelihood, and MLE}{
\[
L(\theta) = \prod_{i=1}^{N} p(x_i;\,\theta), \qquad
\ell(\theta) = \log L(\theta) = \sum_{i=1}^{N} \log p(x_i;\,\theta).
\]
\[
\hat{\theta}_\mathrm{ML} = \arg\max_\theta\; \ell(\theta).
\]
The log transform converts the product into a sum and has the same maximizer since $\log$ is monotone.
}

\thm{MLE for the Gaussian Distribution}{
Given $N$ i.i.d.\ samples from $\mathcal{N}(\mu,\sigma^2)$, the MLEs are:
\[
\mu_\mathrm{ML} = \frac{1}{N}\sum_{i=1}^{N} x_i, \qquad
\sigma^2_\mathrm{ML} = \frac{1}{N}\sum_{i=1}^{N}(x_i - \mu_\mathrm{ML})^2.
\]
}

\pf{
The log-likelihood expands as:
\begin{align*}
\ell(\mu,\sigma^2)
&= \sum_{i=1}^N \log \frac{1}{\sigma\sqrt{2\pi}}\exp\!\left(-\frac{(x_i-\mu)^2}{2\sigma^2}\right) \\
&= -N\log\sigma - \frac{N}{2}\log(2\pi) - \frac{1}{2\sigma^2}\sum_{i=1}^N (x_i - \mu)^2.
\end{align*}

\textbf{Estimating $\mu$:} Set $\pd{\ell}{\mu} = 0$:
\[
\pd{\ell}{\mu} = \frac{1}{\sigma^2}\sum_{i=1}^N (x_i - \mu) = 0 \implies \mu_\mathrm{ML} = \frac{1}{N}\sum_{i=1}^N x_i.
\]

\textbf{Estimating $\sigma^2$:} Let $\tau = \sigma^2$ and set $\pd{\ell}{\tau} = 0$:
\[
\pd{\ell}{\tau} = -\frac{N}{2\tau} + \frac{1}{2\tau^2}\sum_{i=1}^N (x_i - \mu)^2 = 0 \implies \sigma^2_\mathrm{ML} = \frac{1}{N}\sum_{i=1}^N (x_i - \mu_\mathrm{ML})^2.
\]
}

\thm{MLE Variance Estimator is Biased}{
$E\!\left[\sigma^2_\mathrm{ML}\right] = \dfrac{N-1}{N}\,\sigma^2 \neq \sigma^2$. The unbiased estimator is the \textbf{sample variance} $\hat{\sigma}^2 = \dfrac{1}{N-1}\sum_{i=1}^N (x_i - \bar{x})^2$.
}

\pf{
Write $\bar{x} = \mu_\mathrm{ML}$. Use the identity $(x_i - \bar{x}) = (x_i - \mu) - (\bar{x} - \mu)$:
\begin{align*}
\sum_{i=1}^N (x_i - \bar{x})^2
&= \sum_{i=1}^N \bigl[(x_i-\mu) - (\bar{x}-\mu)\bigr]^2 \\
&= \sum_{i=1}^N (x_i-\mu)^2
   - 2(\bar{x}-\mu)\underbrace{\sum_{i=1}^N (x_i-\mu)}_{=\,N(\bar{x}-\mu)}
   + N(\bar{x}-\mu)^2 \\
&= \sum_{i=1}^N (x_i-\mu)^2 - N(\bar{x}-\mu)^2.
\end{align*}
Taking expectations (with $E[(x_i - \mu)^2] = \sigma^2$ and $\mathrm{Var}[\bar{x}] = \sigma^2/N$):
\[
E\!\left[\sigma^2_\mathrm{ML}\right]
= \frac{1}{N}\!\left(N\sigma^2 - N\cdot\frac{\sigma^2}{N}\right)
= \sigma^2 - \frac{\sigma^2}{N}
= \frac{N-1}{N}\,\sigma^2.
\]
}

\thm{Strong Law of Large Numbers}{
Let $X_1, X_2, \ldots$ be i.i.d.\ with $E[X_i] = \mu < \infty$. Then, with probability 1:
\[
\frac{1}{n}\sum_{i=1}^n X_i \;\xrightarrow{\;\mathrm{a.s.}\;}\; \mu \quad \text{as } n\to\infty.
\]
This justifies using the empirical mean as a proxy for the true mean, and more broadly underpins why ERM works.
}

\section{Multiple Random Variables}

\defn{Joint, Marginal, and Conditional Distributions}{
For two random variables $X$ and $Y$:
\begin{align*}
&\textbf{Joint PMF (discrete):} & p_{XY}(x,y) &= P(X=x,\,Y=y) \\
&\textbf{Marginal PMF:} & p_X(x) &= \textstyle\sum_y p_{XY}(x,y) \\
&\textbf{Joint PDF (continuous):} & f_{XY}(x,y) &= \dfrac{\partial^2 F_{XY}(x,y)}{\partial x\,\partial y} \\
&\textbf{Marginal PDF:} & f_X(x) &= \displaystyle\int f_{XY}(x,y)\,dy \\
&\textbf{Conditional PDF:} & f_{Y|X=x}(y) &= \dfrac{f_{XY}(x,y)}{f_X(x)}
\end{align*}
}

\thm{Bayes' Theorem}{
\[
P(B \mid A) = \frac{P(A \mid B)\,P(B)}{P(A)}.
\]
For continuous random variables:
\[
f_{Y|X=x}(y) = \frac{f_{X|Y=y}(x)\,f_Y(y)}{f_X(x)}.
\]
}

\defn{Independence}{
$X$ and $Y$ are \textbf{independent} if $f_{XY}(x,y) = f_X(x)\,f_Y(y)$ (equivalently, $p_{XY}(x,y) = p_X(x)\,p_Y(y)$ for discrete). Independence implies $E[XY] = E[X]E[Y]$.
}

\defn{Covariance and Pearson Correlation}{
\[
\mathrm{Cov}[X,Y] = E\!\bigl[(X - E[X])(Y - E[Y])\bigr] = E[XY] - E[X]E[Y].
\]
\[
\rho_{X,Y} = \frac{\mathrm{Cov}[X,Y]}{\sigma_X\,\sigma_Y} \in [-1,\,1].
\]
Independence implies $\mathrm{Cov}[X,Y] = 0$, but the converse is false in general.
$\rho$ measures linear dependence: $|\rho| = 1$ iff $Y = aX + b$ for some $a \neq 0$.
}

For a random vector $\vect{X} = (X_1,\ldots,X_n)^\top$, the \textbf{covariance matrix} is:
\[
\Sigma = E\!\left[(\vect{X} - E[\vect{X}])(\vect{X} - E[\vect{X}])^\top\right], \quad \Sigma_{ij} = \mathrm{Cov}[X_i, X_j].
\]
$\Sigma$ is always symmetric and positive semidefinite.

The \textbf{sample} counterparts from observations $\{(x_i, y_i)\}_{i=1}^N$:
\[
\bar{x} = \frac{1}{N}\sum_{i=1}^N x_i, \quad
\hat{\sigma}^2 = \frac{1}{N-1}\sum_{i=1}^N (x_i - \bar{x})^2, \quad
\hat{\sigma}^2_{XY} = \frac{1}{N-1}\sum_{i=1}^N (x_i - \bar{x})(y_i - \bar{y}).
\]

\section{Multivariate Gaussian Distribution}

\defn{Multivariate Gaussian}{
A random vector $\vect{X} \in \R^n$ follows $\mathcal{N}(\vect{\mu}, \Sigma)$ with mean $\vect{\mu} \in \R^n$ and PSD covariance $\Sigma \in \R^{n\times n}$ if:
\[
f(\vect{x};\,\vect{\mu},\Sigma) = \frac{1}{(2\pi)^{n/2}|\Sigma|^{1/2}} \exp\!\left(-\frac{1}{2}(\vect{x} - \vect{\mu})^\top \Sigma^{-1}(\vect{x} - \vect{\mu})\right).
\]
The quantity $(\vect{x} - \vect{\mu})^\top \Sigma^{-1}(\vect{x} - \vect{\mu})$ is the squared \textbf{Mahalanobis distance}, which reduces to $\frac{\|\vect{x}-\vect{\mu}\|^2}{\sigma^2}$ when $\Sigma = \sigma^2 I$.
}

% ------------------------------------------------------------------------------
\chapter{Linear Algebra}
% ------------------------------------------------------------------------------

The structure of this chapter follows a hierarchy from general to structured:
\[
\text{General matrices} \;\supset\; \text{Diagonalizable} \;\supset\; \text{Symmetric} \;\supset\; \text{PSD.}
\]
More structure yields nicer properties that ML heavily exploits.

\section{Vectors and Matrices}

We work with real-valued matrices throughout. Vectors are column vectors by default: $\vect{x} \in \R^n$ means $\vect{x} \in \R^{n \times 1}$.

\defn{Basic Operations}{
For $\vect{x}, \vect{y} \in \R^n$ and $A \in \R^{m\times n}$:
\begin{itemize}
    \item \textbf{Inner product:} $\vect{x}^\top \vect{y} = \sum_{i=1}^n x_i y_i \in \R$
    \item \textbf{Outer product:} $\vect{x}\vect{y}^\top \in \R^{n\times n}$, with $(\vect{x}\vect{y}^\top)_{ij} = x_i y_j$
    \item \textbf{Matrix--vector product:} $(A\vect{x})_i = \sum_{j=1}^n A_{ij} x_j$
    \item \textbf{Transpose:} $(A^\top)_{ij} = A_{ji}$;\quad $(AB)^\top = B^\top A^\top$
\end{itemize}
}

\defn{Vector Norms}{
The $\ell_p$ norm of $\vect{x} \in \R^n$: $\norm{\vect{x}}_p = \left(\sum_{i=1}^n |x_i|^p\right)^{1/p}$.
\begin{itemize}
    \item $\norm{\vect{x}}_1 = \sum_i |x_i|$ \quad (sparsity-inducing; used in Lasso)
    \item $\norm{\vect{x}}_2 = \sqrt{\sum_i x_i^2}$ \quad (Euclidean; default $\norm{\vect{x}}$)
    \item $\norm{\vect{x}}_\infty = \max_i |x_i|$
\end{itemize}
}

\section{Matrix Properties}

\thm{Equivalences for Invertible Matrices}{
For $A \in \R^{n\times n}$, the following are equivalent:
\begin{enumerate}
    \item $A$ is invertible ($A^{-1}$ exists with $AA^{-1} = I$).
    \item $\det(A) \neq 0$.
    \item $\mathrm{rank}(A) = n$ (full rank).
    \item $A\vect{x} = \vect{0}$ implies $\vect{x} = \vect{0}$ (trivial null space).
    \item The columns (and rows) of $A$ are linearly independent.
\end{enumerate}
}

\rmkb{The \textbf{trace} $\mathrm{tr}(A) = \sum_i A_{ii}$ equals the sum of eigenvalues; the \textbf{determinant} $\det(A)$ equals their product. Both are invariant under similarity transforms: $\mathrm{tr}(P^{-1}AP) = \mathrm{tr}(A)$.}

\section{Eigendecomposition}

\defn{Eigenvalue and Eigenvector}{
A scalar $\lambda \in \R$ and non-zero vector $\vect{v} \in \R^n$ satisfy
\[
A\vect{v} = \lambda\vect{v}
\]
if $\lambda$ is an \textbf{eigenvalue} and $\vect{v}$ the corresponding \textbf{eigenvector} of $A \in \R^{n\times n}$. The eigenvalues are the roots of the \textbf{characteristic polynomial} $\det(A - \lambda I) = 0$.
}

\thm{Spectral Theorem for Symmetric Matrices}{
If $A \in \R^{n\times n}$ is symmetric ($A = A^\top$), then:
\begin{enumerate}
    \item All eigenvalues are real.
    \item Eigenvectors for distinct eigenvalues are orthogonal.
    \item $A$ is \textbf{orthogonally diagonalizable}: $A = Q\Lambda Q^\top$ where $Q$ is orthogonal ($Q^\top Q = I$), its columns are eigenvectors, and $\Lambda = \mathrm{diag}(\lambda_1, \ldots, \lambda_n)$.
\end{enumerate}
}

\section{Positive Semidefinite Matrices}

\defn{Positive Semidefinite (PSD)}{
A symmetric matrix $A \in \R^{n\times n}$ is \textbf{positive semidefinite} ($A \succeq 0$) if:
\[
\vect{x}^\top A \vect{x} \geq 0 \quad \text{for all } \vect{x} \in \R^n.
\]
It is \textbf{positive definite} ($A \succ 0$) if equality holds only for $\vect{x} = \vect{0}$.
}

\thm{PSD Equivalences}{
For a symmetric $A \in \R^{n\times n}$, the following are equivalent:
\begin{enumerate}
    \item $A \succeq 0$ (by definition).
    \item All eigenvalues $\lambda_i \geq 0$.
    \item $A = B^\top B$ for some matrix $B$.
\end{enumerate}
}

\rmkb{PSD matrices arise as covariance matrices ($\Sigma \succeq 0$ always), and as Hessians of convex functions. The 2nd-order convexity test ($\nabla^2 f \succeq 0$) is precisely a PSD check.}

\section{Singular Value Decomposition}

\thm{Singular Value Decomposition}{
Every matrix $A \in \R^{m\times n}$ of rank $r$ admits:
\[
A = U\Sigma V^\top,
\]
where $U \in \R^{m\times m}$, $V \in \R^{n\times n}$ are orthogonal, and $\Sigma \in \R^{m\times n}$ is diagonal with entries $\sigma_1 \geq \cdots \geq \sigma_r > 0 = \sigma_{r+1} = \cdots$ called \textbf{singular values}.
}

The columns of $U$ are eigenvectors of $AA^\top$; columns of $V$ are eigenvectors of $A^\top A$; $\sigma_i = \sqrt{\lambda_i(A^\top A)}$. SVD is the computational foundation of PCA.

% ------------------------------------------------------------------------------
\chapter{Multivariate Calculus}
% ------------------------------------------------------------------------------

\section{Gradient}

\defn{Gradient}{
For $f: \R^n \to \R$, the \textbf{gradient} is the vector of partial derivatives:
\[
\grad f(\vect{x}) = \nabla f(\vect{x}) = \begin{pmatrix}\pd{f}{x_1} \\ \vdots \\ \pd{f}{x_n}\end{pmatrix} \in \R^n.
\]
$\grad f(\vect{x})$ points in the direction of \emph{steepest ascent}; $-\grad f(\vect{x})$ toward steepest descent.
}

\section{Jacobian and Hessian}

\defn{Jacobian}{
For $\vect{f}: \R^n \to \R^m$, the \textbf{Jacobian} $J \in \R^{m\times n}$ is:
\[
J_{ij} = \pd{f_i}{x_j}, \qquad J = \begin{pmatrix}\grad f_1^\top \\ \vdots \\ \grad f_m^\top\end{pmatrix}.
\]
}

\defn{Hessian}{
For $f: \R^n \to \R$, the \textbf{Hessian} $H = \nabla^2 f(\vect{x}) \in \R^{n\times n}$ is the matrix of second partial derivatives:
\[
H_{ij} = \frac{\partial^2 f}{\partial x_i\,\partial x_j}, \qquad
\nabla^2 f(\vect{x}) = \begin{pmatrix}
\pd{^2 f}{x_1^2} & \cdots & \dfrac{\partial^2 f}{\partial x_1\partial x_n} \\[6pt]
\vdots & \ddots & \vdots \\[4pt]
\dfrac{\partial^2 f}{\partial x_n\partial x_1} & \cdots & \pd{^2 f}{x_n^2}
\end{pmatrix}.
\]
By Schwarz's theorem, $H$ is symmetric whenever mixed partials are continuous.
}

\section{Chain Rule and Matrix Calculus}

\prop{
\textbf{Univariate:} $\dfrac{d}{dt}f(g(t)) = f'(g(t))\cdot g'(t)$.

\textbf{Multivariate:} If $\vect{x} = \vect{g}(\vect{t})$ and $f: \R^n \to \R$:
\[
\pd{f}{t_k} = \sum_{i=1}^n \pd{f}{x_i}\,\pd{x_i}{t_k}, \qquad
\nabla_{\vect{t}} f = J_{\vect{g}}^\top\, \nabla_{\vect{x}} f.
\]
}

\fact{
Useful matrix calculus identities (where $A$ is a constant matrix):
\begin{align*}
\nabla_{\vect{x}}\,(A\vect{x}) &= A^\top \\
\nabla_{\vect{x}}\,(\vect{x}^\top A\vect{x}) &= (A + A^\top)\vect{x} = 2A\vect{x} \quad \text{(if $A$ symmetric)}
\end{align*}
The second identity is used directly in the normal equations for linear regression.
}

% ------------------------------------------------------------------------------
\chapter{Convex Optimization}
\label{ch:convex}
% ------------------------------------------------------------------------------

\section{Convex Sets and Functions}

\defn{Convex Set}{
$\mathcal{C} \subseteq \R^d$ is \textbf{convex} if for any $\vect{x}, \vect{y} \in \mathcal{C}$ and $\theta \in [0,1]$:
\[
\theta\vect{x} + (1-\theta)\vect{y} \in \mathcal{C}.
\]
Examples: $\R^d$; affine subspaces $\{\vect{x} \mid A\vect{x} = \vect{b}\}$; polyhedra $\{\vect{x} \mid A\vect{x} \leq \vect{b}\}$. Intersections of convex sets are convex.
}

\defn{Convex Function}{
$f: \mathcal{D} \to \R$ (with convex $\mathcal{D}$) is \textbf{convex} if for all $\vect{x}_1, \vect{x}_2 \in \mathcal{D}$ and $\lambda \in [0,1]$:
\[
f\!\left(\lambda\vect{x}_1 + (1-\lambda)\vect{x}_2\right) \leq \lambda f(\vect{x}_1) + (1-\lambda)f(\vect{x}_2).
\]
$f$ is \textbf{strictly convex} if strict inequality holds for $\vect{x}_1 \neq \vect{x}_2$, $\lambda \in (0,1)$. $f$ is \textbf{concave} if $-f$ is convex.
}

\thm{Global Optimality of Critical Points}{
If $f$ is convex and $\vect{x}^*$ satisfies $\nabla f(\vect{x}^*) = \vect{0}$ (or the gradient does not exist there), then $\vect{x}^*$ is a \textbf{global minimum} of $f$.
}

This is the central reason ML prefers convex losses: any optimizer that finds a critical point has found the global solution.

\section{Verifying Convexity}

\thm{Three Convexity Tests --- Scalar}{
For $f: \R \to \R$, each of the following is equivalent to convexity:
\begin{enumerate}
    \item \textbf{0th-order:} $f(\lambda x_1 + (1-\lambda)x_2) \leq \lambda f(x_1) + (1-\lambda)f(x_2)$ \quad (definition).
    \item \textbf{1st-order (tangent line below):} $f(x_2) \geq f(x_1) + f'(x_1)(x_2 - x_1)$ for all $x_1, x_2$.
    \item \textbf{2nd-order (non-negative curvature):} $f''(x) \geq 0$ for all $x$.
\end{enumerate}
}

\thm{Three Convexity Tests --- Multivariate}{
For $f: \R^n \to \R$ with convex domain, each of the following is equivalent to convexity:
\begin{enumerate}
    \item \textbf{0th-order:} $f(\lambda\vect{x} + (1-\lambda)\vect{y}) \leq \lambda f(\vect{x}) + (1-\lambda)f(\vect{y})$.
    \item \textbf{1st-order (tangent hyperplane below):} $f(\vect{y}) \geq f(\vect{x}) + (\nabla f(\vect{x}))^\top(\vect{y} - \vect{x})$ for all $\vect{x}, \vect{y}$.
    \item \textbf{2nd-order (Hessian PSD):} $\nabla^2_{\vect{x}} f(\vect{x}) \succeq 0$ for all $\vect{x}$.
\end{enumerate}
Equivalently: $f$ is convex iff its restriction to every line is convex.
}

\ex{
$f(\vect{x}) = 2x_1^2 + x_2^2$ has $H = \begin{pmatrix}4&0\\0&2\end{pmatrix}$. Eigenvalues $4, 2 > 0$, so $H \succ 0$ and $f$ is strictly convex.

$f(\vect{x}) = 2x_1^2 - x_2^2$ has $H = \begin{pmatrix}4&0\\0&-2\end{pmatrix}$. Eigenvalue $-2 < 0$, so $H \not\succeq 0$ and $f$ is not convex.

$f(\vect{x}) = x_1 + 2x_1^2 + x_2^2 - x_1 x_2 - 2$ has $H = \begin{pmatrix}4&-1\\-1&2\end{pmatrix}$. Eigenvalues $3 \pm \sqrt{2} > 0$, so $f$ is strictly convex.
}

\section{Gradient Descent}

\defn{Gradient Descent}{
To minimize $f(\vect{w})$, initialize $\vect{w}_0$ and iterate:
\[
\vect{w}_{t+1} = \vect{w}_t - \eta\,\nabla f(\vect{w}_t), \quad t = 0, 1, 2, \ldots
\]
until $\|\nabla f(\vect{w}_t)\| < \varepsilon$ or max iterations reached. The scalar $\eta > 0$ is the \textbf{learning rate}.
}

\rmkb{
\textbf{Choosing $\eta$}: too small $\Rightarrow$ slow convergence; too large $\Rightarrow$ oscillation or divergence. For an $L$-smooth function (Lipschitz gradient constant $L$), the theoretically safe step size is $\eta = 1/L$.
}

\thm{GD Convergence for Convex Functions}{
Let $f$ be convex and $L$-smooth ($\|\nabla f(\vect{x}) - \nabla f(\vect{y})\| \leq L\|\vect{x}-\vect{y}\|$). With step size $\eta = 1/L$ and global minimizer $\vect{w}^*$:
\[
f(\vect{w}_t) - f(\vect{w}^*) \leq \frac{L\|\vect{w}_0 - \vect{w}^*\|^2}{2t} = O\!\left(\frac{1}{t}\right).
\]
For $\mu$-strongly convex $f$, the rate improves to $O\!\left(\left(1 - \frac{\mu}{L}\right)^t\right)$ (linear/geometric convergence).
}

\section{Stochastic Gradient Descent}

When $f(\vect{w}) = \frac{1}{N}\sum_{i=1}^N \ell_i(\vect{w})$ (ERM), computing $\nabla f$ requires $O(N)$ evaluations per step.

\defn{SGD and Mini-Batch SGD}{
\textbf{SGD} samples a single index $i_t \sim \mathrm{Uniform}\{1,\ldots,N\}$ per step:
\[
\vect{w}_{t+1} = \vect{w}_t - \eta\,\nabla\ell_{i_t}(\vect{w}_t).
\]
\textbf{Mini-batch SGD} averages over a batch $\mathcal{B} \subset \{1,\ldots,N\}$ of size $B$:
\[
\vect{w}_{t+1} = \vect{w}_t - \frac{\eta}{B}\sum_{i \in \mathcal{B}} \nabla\ell_i(\vect{w}_t).
\]
}

The key property: $E_{i_t}[\nabla\ell_{i_t}(\vect{w})] = \nabla f(\vect{w})$. SGD uses an \textbf{unbiased estimator} of the true gradient, enabling progress while spending only $O(B)$ per update.

% ------------------------------------------------------------------------------
\chapter{Model Evaluation}
\label{ch:eval}
% ------------------------------------------------------------------------------

\section{The Bias-Variance Decomposition}

Suppose the true relationship is $y = f(\vect{x}) + \varepsilon$ with $\varepsilon \sim \mathcal{N}(0, \sigma^2_\varepsilon)$. Let $\hat{f}$ be a model trained on a random dataset $\mathcal{D}$.

\thm{Bias-Variance-Noise Decomposition}{
The expected squared error at a fixed point $\vect{x}$ decomposes as:
\[
E_{\mathcal{D},\varepsilon}\!\left[(y - \hat{f}(\vect{x}))^2\right]
= \underbrace{\Bigl(f(\vect{x}) - E_\mathcal{D}[\hat{f}(\vect{x})]\Bigr)^2}_{\mathrm{Bias}^2}
+ \underbrace{E_\mathcal{D}\!\left[\Bigl(\hat{f}(\vect{x}) - E_\mathcal{D}[\hat{f}(\vect{x})]\Bigr)^2\right]}_{\mathrm{Variance}}
+ \underbrace{\sigma^2_\varepsilon}_{\mathrm{Noise}}.
\]
}

\pf{
Let $\bar{f} = E_\mathcal{D}[\hat{f}(\vect{x})]$. First isolate the noise:
\begin{align*}
E\!\left[(y - \hat{f})^2\right]
&= E\!\left[(f + \varepsilon - \hat{f})^2\right] \\
&= E\!\left[(f - \hat{f})^2\right] + 2\underbrace{E[\varepsilon]}_{=\,0}\,E[f - \hat{f}] + E[\varepsilon^2] \\
&= E\!\left[(f - \hat{f})^2\right] + \sigma^2_\varepsilon.
\end{align*}
Now add and subtract $\bar{f}$:
\begin{align*}
E\!\left[(f - \hat{f})^2\right]
&= E\!\left[(f - \bar{f} + \bar{f} - \hat{f})^2\right] \\
&= (f - \bar{f})^2 + 2(f - \bar{f})\underbrace{E[\bar{f} - \hat{f}]}_{=\,0} + E\!\left[(\bar{f} - \hat{f})^2\right] \\
&= \mathrm{Bias}^2[\hat{f}] + \mathrm{Var}[\hat{f}]. \qquad\square
\end{align*}
}

The decomposition reveals a fundamental tradeoff:
\begin{itemize}
    \item \textbf{High bias, low variance} (underfitting): model too simple; increase capacity.
    \item \textbf{Low bias, high variance} (overfitting): model too complex; regularize or collect more data.
    \item \textbf{Noise $\sigma^2_\varepsilon$}: irreducible; no model can eliminate it.
\end{itemize}

\section{Cross-Validation}

\defn{$k$-Fold Cross-Validation}{
Partition training data into $k$ equal folds. For each fold $j = 1,\ldots,k$: train on the remaining $k-1$ folds, evaluate on fold $j$. Report the average metric over all $k$ folds.
}

$k$-fold CV provides a lower-variance estimate of $R(h)$ than a single split, especially for small $N$. Common choices: $k = 5$ or $k = 10$. Special case $k = N$: \textbf{leave-one-out CV (LOOCV)}.

\section{Classification Metrics}

For binary classification with predictions $\hat{y} \in \{0,1\}$ and true labels $y \in \{0,1\}$:

\begin{center}
\renewcommand{\arraystretch}{1.3}
\begin{tabular}{l|cc}
& Predicted Positive & Predicted Negative \\
\hline
Actual Positive & TP & FN \\
Actual Negative & FP & TN \\
\end{tabular}
\end{center}

\defn{Precision, Recall, and F1}{
\[
\text{Accuracy} = \frac{TP + TN}{TP + FP + TN + FN}, \qquad
\text{Precision} = \frac{TP}{TP + FP}, \qquad
\text{Recall} = \frac{TP}{TP + FN}.
\]
\[
\text{F1} = \frac{2 \cdot \text{Precision} \cdot \text{Recall}}{\text{Precision} + \text{Recall}}.
\]
}

\rmkb{
Accuracy is misleading on imbalanced datasets. A classifier that predicts ``always negative'' on a 99\%-negative dataset gets 99\% accuracy while detecting nothing. Use Precision, Recall, or F1.

The \textbf{ROC curve} plots True Positive Rate (Recall) vs.\ False Positive Rate $\frac{FP}{FP+TN}$ as the decision threshold varies. The \textbf{AUC} (area under the ROC curve) summarizes performance across all thresholds; $\text{AUC} = 1$ is perfect, $\text{AUC} = 0.5$ is random.
}

% ==============================================================================
% PART II: SUPERVISED LEARNING
% ==============================================================================

\part{Supervised Learning}

% ------------------------------------------------------------------------------
\chapter{Linear Regression}
% ------------------------------------------------------------------------------

\section{The Linear Model}

\defn{Linear Regression}{
Given input $\vect{x} = (x_1, \ldots, x_d)^T \in \R^d$, the linear model predicts a real-valued output via
\[
f(\vect{x}; \vect{w}) = w_0 + \sum_{j=1}^{d} w_j x_j = \vect{w}^T \vect{x},
\]
where the bias $w_0$ is absorbed by prepending a constant 1 to $\vect{x}$, so $\vect{w} \in \R^{d+1}$.
}

We stack $n$ training examples into a design matrix $X \in \R^{n \times (d+1)}$ (rows are data points) and label vector $\vect{y} \in \R^n$. The predicted output vector is $\hat{\vect{y}} = X\vect{w}$.

\section{Loss Functions}

\defn{Mean Squared Error}{
\[
L_{\mathrm{MSE}}(\vect{w}) = \frac{1}{n}\sum_{i=1}^n \bigl(y_i - \vect{w}^T\vect{x}_i\bigr)^2 = \frac{1}{n}\norm{X\vect{w} - \vect{y}}^2.
\]
MSE penalizes large errors quadratically and is differentiable everywhere. It is sensitive to outliers.
}

\defn{Mean Absolute Error}{
\[
L_{\mathrm{MAE}}(\vect{w}) = \frac{1}{n}\sum_{i=1}^n |y_i - \vect{w}^T \vect{x}_i|.
\]
Robust to outliers but not differentiable at zero; requires subgradient methods.
}

\defn{Huber Loss}{
A smooth compromise between MSE and MAE, controlled by $\delta > 0$:
\[
H_\delta(r) = \begin{cases} \tfrac{1}{2}r^2 & |r| \leq \delta, \\ \delta\bigl(|r| - \tfrac{1}{2}\delta\bigr) & |r| > \delta. \end{cases}
\]
Quadratic near zero (differentiable), linear in the tails (robust).
}

\rmkb{The loss choice encodes a noise assumption. MSE corresponds to Gaussian noise $\epsilon \sim \mathcal{N}(0, \sigma^2)$; MAE to Laplacian noise; Huber to a robust mixture. For high-leverage points or outliers, MAE or Huber is preferred.}

\section{Ordinary Least Squares}

\thm{Normal Equations}{
The unique minimizer of $L_{\mathrm{MSE}}$ satisfies
\[
X^T X\, \hat{\vect{w}} = X^T \vect{y},
\]
and when $X$ has full column rank, $\hat{\vect{w}} = (X^T X)^{-1} X^T \vect{y}$.
}

\pf{
Expanding $L(\vect{w}) = \frac{1}{n}(X\vect{w}-\vect{y})^T(X\vect{w}-\vect{y})$ and differentiating:
\[
\grad_{\vect{w}} L = \frac{2}{n} X^T(X\vect{w} - \vect{y}) = \vect{0}
\quad\Longrightarrow\quad
X^T X \vect{w} = X^T \vect{y}.
\]
The matrix $X^T X$ is positive semi-definite. Full column rank of $X$ implies $X^T X \succ 0$, giving a unique solution via matrix inversion.
}

\fact{The level set $\{\vect{w} \mid (\vect{w}-\hat{\vect{w}})^T X^T X (\vect{w}-\hat{\vect{w}}) = C\}$ is an ellipsoid whose axis lengths scale with the eigenvalues of $(X^T X)^{-1}$. The gradient descent trajectory descends perpendicular to these contours.}

Gradient descent iterates $\vect{w} \leftarrow \vect{w} - \eta\,\grad L$, which expands to the per-weight update:
\[
w_j \leftarrow w_j + \frac{2\eta}{n}\sum_{i=1}^n \bigl[y_i - f(\vect{x}_i; \vect{w})\bigr]\, x_{i,j}.
\]

\section{Probabilistic Interpretation}
\label{sec:linreg_mle}

Assume $y_i = \vect{w}^T\vect{x}_i + \epsilon_i$ with $\epsilon_i \overset{\mathrm{iid}}{\sim} \mathcal{N}(0, \sigma^2)$. Then
\[
P(y_i \mid \vect{x}_i; \vect{w}) = \frac{1}{\sqrt{2\pi}\sigma}\exp\!\left(-\frac{(y_i - \vect{w}^T\vect{x}_i)^2}{2\sigma^2}\right).
\]

\prop{MLE $\equiv$ OLS under Gaussian Noise}{
\[
\hat{\vect{w}}_{\mathrm{MLE}}
= \arg\max_{\vect{w}}\sum_{i=1}^n \log P(y_i|\vect{x}_i;\vect{w})
= \arg\min_{\vect{w}} \frac{1}{n}\sum_{i=1}^n (y_i - \vect{w}^T\vect{x}_i)^2
= \hat{\vect{w}}_{\mathrm{OLS}}.
\]
The negative log-likelihood of i.i.d.\ Gaussian noise is exactly the MSE.
}

\section{Collinearity}

\defn{Collinearity}{
Two predictors $x_j$ and $x_k$ are \textbf{collinear} (or multicollinear in a group) if they are approximately linearly dependent. Perfectly collinear predictors make $X^T X$ singular.
}

\fact{When predictors are highly correlated, the MSE contours form a very elongated ellipsoid: the eigenvalues of $X^T X$ differ by orders of magnitude. The OLS estimate $\hat{\vect{w}} = (X^T X)^{-1} X^T \vect{y}$ is extremely sensitive to small perturbations in the data. Ridge regression (Chapter~\ref{ch:regularization}) resolves this by replacing $X^T X$ with $X^T X + \lambda I$, which is always invertible.}

\ex{Suppose $x_2 = 2x_1$ exactly. Then the $2 \times 2$ matrix $X^T X$ for $(x_1, x_2)$ has a zero eigenvalue, and infinitely many $\hat{\vect{w}}$ achieve the same MSE. Any combination $w_1 + 2w_2 = \text{const}$ lies in the null space.}

% ------------------------------------------------------------------------------
\chapter{Logistic Regression}
% ------------------------------------------------------------------------------

\section{From Regression to Classification}

For binary classification $y \in \{0, 1\}$, we want to model $P(Y=1 \mid X=\vect{x})$. Applying a linear model directly is problematic: outputs can fall outside $[0,1]$. We instead compose the linear function with the \textbf{sigmoid}:

\defn{Sigmoid (Logistic) Function}{
\[
\sigma(z) = \frac{1}{1+e^{-z}} = \frac{e^z}{e^z + 1}, \quad z \in \R,\quad \sigma(z) \in (0,1).
\]
Properties: $\sigma'(z) = \sigma(z)(1-\sigma(z))$; symmetric about 0; $\sigma(z) \to 1$ as $z \to +\infty$.
}

The logistic regression model is:
\[
p(\vect{x}) = P(Y=1 \mid X=\vect{x}) = \sigma(\vect{w}^T\vect{x}) = \frac{1}{1+e^{-\vect{w}^T\vect{x}}}.
\]
The decision rule is $\hat{y} = \mathbf{1}[p(\vect{x}) > 0.5]$, equivalently $\hat{y} = \mathbf{1}[\vect{w}^T\vect{x} > 0]$. The \textbf{decision boundary} is the hyperplane $\vect{w}^T\vect{x} = 0$.

\section{Log-Odds and the Logit Link}

\prop{Log-Odds Interpretation}{
Rearranging the logistic model gives the \textbf{logit} (log-odds) function:
\[
\log\frac{p(\vect{x})}{1-p(\vect{x})} = \vect{w}^T\vect{x}.
\]
Logistic regression is a \emph{linear model for the log-odds}. Each weight $w_j$ shifts the log-odds by $w_j$ per unit increase in $x_j$.
}

\section{Learning via Maximum Likelihood}

\defn{Cross-Entropy Loss}{
Since $P(Y=y\mid X=\vect{x}) = p(\vect{x})^y [1-p(\vect{x})]^{1-y}$, the negative log-likelihood over the training set is the \textbf{binary cross-entropy}:
\[
L(\vect{w}) = -\sum_{i=1}^n \bigl\{y_i \log p(\vect{x}_i; \vect{w}) + (1-y_i)\log(1-p(\vect{x}_i;\vect{w}))\bigr\}.
\]
}

\prop{Convexity of Cross-Entropy}{
The cross-entropy loss is convex in $\vect{w}$, so gradient descent converges to the global minimum (if one exists). In contrast, using MSE on $p(\vect{x}) = \sigma(\vect{w}^T\vect{x})$ yields a non-convex objective with many local minima.
}

\pf{(Sketch) The Hessian of $L(\vect{w})$ equals $X^T D X$ where $D = \mathrm{diag}(p_i(1-p_i))$ is a diagonal matrix with non-negative entries. Hence the Hessian is positive semi-definite everywhere, confirming convexity.}

\section{Gradient and Gradient Descent}

The gradient of the cross-entropy loss is remarkably clean:
\[
\grad_{\vect{w}} L = -\sum_{i=1}^n \bigl[y_i - p(\vect{x}_i; \vect{w})\bigr]\,\vect{x}_i.
\]
The update rule for each weight $w_j$ is:
\[
w_j \leftarrow w_j - \eta \sum_{i=1}^n -\bigl[y_i - p(\vect{x}_i; \vect{w})\bigr]\,x_{i,j}.
\]

\rmkb{The gradient has the same form as in linear regression, with the prediction $p(\vect{x}_i;\vect{w}) = \sigma(\vect{w}^T\vect{x}_i)$ replacing $\vect{w}^T\vect{x}_i$. This structural similarity makes both models easy to unify in a single codebase.}

\prop{Well-Calibration}{
At the gradient optimum, setting $\grad L = 0$ and examining the intercept term yields $\sum_i y_i = \sum_i p(\vect{x}_i; \hat{\vect{w}})$: the model's average predicted probability equals the empirical class frequency. Logistic regression is \emph{well-calibrated} in this sense.
}

\section{Multiclass: Softmax Regression}

For $K$ classes, we learn weight vectors $\vect{w}_1, \ldots, \vect{w}_K$ and compute class logits $z_k = w_{k0} + \vect{w}_k^T\vect{x}$. The \textbf{softmax} function converts logits to a probability distribution:

\defn{Softmax}{
\[
P(Y=k \mid X=\vect{x}) = \frac{\exp(z_k)}{\sum_{k'=1}^{K}\exp(z_{k'})}, \quad k = 1, \ldots, K.
\]
The outputs are non-negative and sum to 1. As $\alpha \to \infty$, $\mathrm{softmax}(\alpha z)$ concentrates all mass on $\arg\max_k z_k$.
}

\fact{Softmax generalizes sigmoid: for $K=2$, $P(Y=1) = \exp(z_1)/(\exp(z_1)+\exp(z_2)) = \sigma(z_1 - z_2)$.}

The loss is the multiclass cross-entropy:
\[
L = -\sum_{i=1}^n \sum_{k=1}^K y_{ik}\log\hat{y}_{ik},
\]
where $\vect{y}_i$ is a one-hot vector and $\hat{y}_{ik} = P(Y=k\mid\vect{x}_i)$.

\section{Limitations}

Logistic regression produces a \emph{linear} decision boundary $\vect{w}^T\vect{x} = 0$. For problems like XOR (where no hyperplane separates the classes), it fails. The natural remedy is to either: (a) engineer nonlinear features $\phi(\vect{x})$, or (b) cascade multiple logistic regression models --- which gives a multilayer neural network (Part~IV).

% ------------------------------------------------------------------------------
\chapter{Regularization}
\label{ch:regularization}
% ------------------------------------------------------------------------------

\section{The Overfitting Problem}

Without constraints, empirical risk minimization can fit training data perfectly while failing to generalize. The bias-variance decomposition (Chapter~\ref{ch:learnproblem}) shows that model complexity is the lever: increasing it reduces bias but raises variance. Regularization adds an explicit complexity penalty to the loss.

\section{Ridge Regression ($\ell_2$ Regularization)}

\defn{Ridge Regression}{
Ridge regression augments the MSE objective with an $\ell_2$ penalty on the weights:
\[
L_{\mathrm{ridge}}(\vect{w}) = \frac{1}{n}\norm{X\vect{w} - \vect{y}}^2 + \lambda \norm{\vect{w}}^2, \quad \lambda > 0.
\]
The bias term $w_0$ is conventionally excluded from penalization.
}

\thm{Ridge Closed Form}{
The unique minimizer is
\[
\hat{\vect{w}}_{\mathrm{ridge}} = (X^T X + \lambda I)^{-1} X^T \vect{y}.
\]
The matrix $X^T X + \lambda I$ is positive definite for any $\lambda > 0$, resolving collinearity.
}

\pf{
$\grad L_{\mathrm{ridge}} = \frac{2}{n}X^T(X\vect{w}-\vect{y}) + 2\lambda\vect{w} = 0$ gives $(X^T X + n\lambda I)\vect{w} = X^T\vect{y}$, and absorbing the factor of $n$ into $\lambda$.
}

\prop{Bayesian Interpretation of Ridge}{
Ridge is equivalent to MAP estimation with a Gaussian prior $\vect{w} \sim \mathcal{N}(\vect{0}, \frac{1}{2\lambda}I)$ on the weights. The $\ell_2$ penalty is the negative log-prior.
}

\rmkb{Geometrically, the Ridge constraint $\norm{\vect{w}}^2 \leq t$ defines a sphere. The solution is where the OLS contours first touch the sphere. All weights are shrunk toward zero, but none are exactly zero.}

\section{Lasso ($\ell_1$ Regularization)}

\defn{Lasso}{
\[
L_{\mathrm{lasso}}(\vect{w}) = \frac{1}{n}\norm{X\vect{w}-\vect{y}}^2 + \lambda \norm{\vect{w}}_1, \quad \norm{\vect{w}}_1 = \sum_{j=1}^d |w_j|.
\]
}

\prop{Sparsity of Lasso}{
The $\ell_1$ constraint region is a cross-polytope (diamond in 2D). Its corners lie on coordinate axes, so the OLS contours frequently touch the constraint boundary at a corner --- setting some weights to \emph{exactly zero}. Lasso performs automatic feature selection.
}

Unlike Ridge, Lasso has no closed-form solution (the $\ell_1$ norm is non-differentiable at 0). It is solved via coordinate descent or proximal gradient methods.

\prop{Bayesian Interpretation of Lasso}{
Lasso is equivalent to MAP estimation with a Laplace prior $p(w_j) \propto \exp(-\lambda|w_j|)$. The heavier tails of the Laplace distribution encourage sparsity.
}

\section{Elastic Net}

\defn{Elastic Net}{
Elastic Net combines $\ell_1$ and $\ell_2$ penalties:
\[
L_{\mathrm{EN}}(\vect{w}) = \frac{1}{n}\norm{X\vect{w}-\vect{y}}^2 + \lambda_1 \norm{\vect{w}}_1 + \lambda_2 \norm{\vect{w}}^2.
\]
It inherits sparsity from Lasso and the grouping effect (correlated features get similar weights) from Ridge.
}

\section{Regularization for Logistic Regression}

The same penalties apply to classification. The regularized cross-entropy objectives are:
\begin{align*}
L_{\mathrm{L2}}(\vect{w}) &= -\sum_{i=1}^n \{y_i \log p_i + (1-y_i)\log(1-p_i)\} + \lambda\sum_{j=1}^d w_j^2, \\
L_{\mathrm{L1}}(\vect{w}) &= -\sum_{i=1}^n \{y_i \log p_i + (1-y_i)\log(1-p_i)\} + \lambda\sum_{j=1}^d |w_j|.
\end{align*}
$\ell_2$ regularization prevents the weights from growing arbitrarily large when classes are linearly separable (which would push logistic outputs to 0/1 and cause numerical issues).

\section{Effect on Bias and Variance}

\prop{Bias-Variance Tradeoff in Ridge}{
As $\lambda \to 0$, $\hat{\vect{w}}_{\mathrm{ridge}} \to \hat{\vect{w}}_{\mathrm{OLS}}$ (zero bias, full variance). As $\lambda \to \infty$, $\hat{\vect{w}}_{\mathrm{ridge}} \to \vect{0}$ (high bias, zero variance). The optimal $\lambda$ is chosen by cross-validation.
}

% ------------------------------------------------------------------------------
\chapter{Support Vector Machines}
% ------------------------------------------------------------------------------

\section{Hard-Margin SVM}

For binary classification with labels $y_i \in \{+1, -1\}$ and a linearly separable dataset, we seek the hyperplane $\vect{w}^T\vect{x} + w_0 = 0$ with \emph{maximum margin}.

\defn{Margin}{
The \textbf{functional margin} of point $(\vect{x}_i, y_i)$ is $y_i(\vect{w}^T\vect{x}_i + w_0)$. The \textbf{geometric margin} (distance from the hyperplane) is $y_i(\vect{w}^T\vect{x}_i + w_0)/\norm{\vect{w}}$. The margin of the classifier is $\min_i y_i(\vect{w}^T\vect{x}_i + w_0)/\norm{\vect{w}}$.
}

By rescaling $\vect{w}$ so that $\min_i y_i(\vect{w}^T\vect{x}_i + w_0) = 1$, the margin becomes $1/\norm{\vect{w}}$. Maximizing the margin is equivalent to:

\defn{Hard-Margin SVM (Primal)}{
\[
\min_{\vect{w}, w_0} \frac{1}{2}\norm{\vect{w}}^2 \quad \text{s.t.} \quad y_i(\vect{w}^T\vect{x}_i + w_0) \geq 1 \;\forall\, i.
\]
The points achieving equality, $y_i(\vect{w}^T\vect{x}_i + w_0) = 1$, are the \textbf{support vectors}.
}

\section{Dual Formulation}

Introducing Lagrange multipliers $\alpha_i \geq 0$ for each constraint gives the dual problem:

\defn{SVM Dual}{
\[
\max_{\boldsymbol{\alpha}} \sum_{i=1}^n \alpha_i - \frac{1}{2}\sum_{i,j} y_i y_j \alpha_i \alpha_j \langle\vect{x}_i, \vect{x}_j\rangle
\quad\text{s.t.}\quad \alpha_i \geq 0,\; \sum_{i=1}^n \alpha_i y_i = 0.
\]
}

\prop{KKT Complementary Slackness}{
At optimality: $\alpha_i [y_i(\vect{w}^T\vect{x}_i + w_0) - 1] = 0$. Thus $\alpha_i > 0$ only for support vectors. The weight vector $\vect{w} = \sum_{i=1}^n \alpha_i y_i \vect{x}_i$ is a linear combination of support vectors alone.
}

Prediction uses only inner products with support vectors:
\[
f(\vect{x}) = w_0 + \sum_{(\vect{x}', y') \in \mathrm{SV}} y' \alpha' \langle \vect{x}', \vect{x} \rangle.
\]

\section{Soft-Margin SVM and Hinge Loss}

When the data is not linearly separable, we allow margin violations via \textbf{slack variables} $\xi_i \geq 0$:

\defn{Soft-Margin SVM}{
\[
\min_{\vect{w}, w_0, \boldsymbol{\xi}} \frac{1}{2}\norm{\vect{w}}^2 + C\sum_{i=1}^n \xi_i
\quad\text{s.t.}\quad y_i(\vect{w}^T\vect{x}_i + w_0) \geq 1 - \xi_i,\; \xi_i \geq 0.
\]
$C > 0$ trades margin width against training error. Substituting $\xi_i = \max(1 - y_i(\vect{w}^T\vect{x}_i + w_0), 0)$ gives the unconstrained objective:
}

\prop{SVM as Regularized Hinge Loss}{
\[
\min_{\vect{w}, w_0}\; \frac{1}{2}\norm{\vect{w}}^2 + C\sum_{i=1}^n \max\!\bigl(1 - y_i(w_0 + \vect{w}^T\vect{x}_i),\, 0\bigr).
\]
The \textbf{hinge loss} $\ell_{\mathrm{hinge}}(y, f) = \max(1-yf, 0)$ is convex and zero when the point is correctly classified with margin $\geq 1$.
}

\section{The Kernel Trick}

The dual SVM objective depends on the data only through inner products $\langle\vect{x}_i, \vect{x}_j\rangle$. If we apply a feature map $\phi: \R^d \to \R^m$ (possibly $m = \infty$), we can replace inner products with a \textbf{kernel function} $k(\vect{x}_i, \vect{x}_j) = \langle\phi(\vect{x}_i), \phi(\vect{x}_j)\rangle$, \emph{without ever computing $\phi(\vect{x})$ explicitly}.

\defn{Common Kernels}{
\begin{center}
\begin{tabular}{lll}
\toprule
Kernel & Formula & Remark \\
\midrule
Linear & $k(\vect{u},\vect{v}) = \vect{u}^T\vect{v}$ & Standard SVM \\
Polynomial & $k(\vect{u},\vect{v}) = (1 + \vect{u}^T\vect{v})^D$ & Degree-$D$ polynomial map \\
RBF / Gaussian & $k(\vect{u},\vect{v}) = \exp(-\gamma\norm{\vect{u}-\vect{v}}^2)$ & Infinite-dimensional $\phi$ \\
\bottomrule
\end{tabular}
\end{center}
}

\prop{Mercer's Condition}{
A symmetric function $k$ is a valid kernel (i.e., corresponds to some $\phi$) if and only if the Gram matrix $K_{ij} = k(\vect{x}_i, \vect{x}_j)$ is positive semi-definite for all data sets.
}

With a kernel, prediction becomes:
\[
f(\vect{x}) = w_0 + \sum_{(\vect{x}',y') \in \mathrm{SV}} \alpha' y'\, k(\vect{x}', \vect{x}).
\]
The RBF kernel is equivalent to an infinite Taylor expansion: $k(\vect{x}_i,\vect{x}_j) = C\sum_{n=0}^\infty \frac{\langle\vect{x}_i,\vect{x}_j\rangle^n}{n!}$, giving it universal approximation capability.

% ------------------------------------------------------------------------------
\chapter{\textit{K}-Nearest Neighbors}
% ------------------------------------------------------------------------------

\section{The Algorithm}

\defn{$K$-Nearest Neighbors (KNN)}{
Given a test point $\vect{x}$, find the $K$ training points $\mathcal{N}_K(\vect{x}) = \{\vect{x}_{(1)}, \ldots, \vect{x}_{(K)}\}$ closest under a distance metric $d(\cdot,\cdot)$.
\begin{itemize}
    \item \textbf{Classification:} $\hat{y} = \mathrm{mode}\{y_i : \vect{x}_i \in \mathcal{N}_K(\vect{x})\}$ (majority vote).
    \item \textbf{Regression:} $\hat{y} = \frac{1}{K}\sum_{\vect{x}_i \in \mathcal{N}_K(\vect{x})} y_i$ (average).
\end{itemize}
KNN is a \textbf{lazy learner}: there is no training phase; all computation occurs at prediction time.
}

\section{Distance Metrics}

The choice of distance metric critically affects KNN performance.

\defn{Minkowski Distance}{
For $p \geq 1$:
\[
d_p(\vect{u}, \vect{v}) = \left(\sum_{j=1}^d |u_j - v_j|^p\right)^{1/p}.
\]
\begin{itemize}
    \item $p=2$: Euclidean distance (most common).
    \item $p=1$: Manhattan (taxicab) distance. More robust to outliers.
    \item $p \to \infty$: Chebyshev distance $\max_j |u_j - v_j|$.
\end{itemize}
}

\rmkb{Features should be normalized before applying KNN. If one feature has range $[0, 1000]$ and another $[0, 1]$, the first will dominate the distance computation.}

\section{Choosing $K$ and the Bias-Variance Tradeoff}

\prop{Effect of $K$}{
Small $K$ (e.g., $K=1$): the decision boundary follows every training point, yielding low bias but high variance (overfitting). Large $K$: smoother boundary, higher bias, lower variance. The optimal $K$ is selected by cross-validation.
}

\prop{Asymptotic Optimality}{
As $n \to \infty$ with $K \to \infty$ and $K/n \to 0$, the 1-NN error rate converges to at most $2R^*$ where $R^*$ is the Bayes error rate. For $K \geq 1$, KNN is universally consistent.
}

\section{Curse of Dimensionality}

\fact{In $d$ dimensions, to capture a fraction $f$ of the data volume in a neighborhood, the side length of the neighborhood must be $f^{1/d}$. For $f = 0.01$ and $d = 10$, this is $0.01^{0.1} \approx 0.63$ --- nearly the full feature range. Nearest neighbors cease to be meaningfully ``near'' in high dimensions.}

\prop{Computational Complexity}{
Naïve KNN requires $O(nd)$ distance computations per query. With $n$ training points and $q$ test points, total cost is $O(qnd)$. Data structures such as $k$-d trees reduce average-case query time to $O(d \log n)$ for low $d$, but lose their advantage as $d$ grows.
}

% ------------------------------------------------------------------------------
\chapter{Na\"{i}ve Bayes}
% ------------------------------------------------------------------------------

\section{Generative vs.\ Discriminative Classifiers}

\defn{Discriminative Model}{
Directly models the posterior $P(Y \mid X)$. Examples: logistic regression, SVM.
}

\defn{Generative Model}{
Models the joint distribution $P(X, Y) = P(X \mid Y)P(Y)$, then applies Bayes' rule to obtain $P(Y \mid X)$.
}

\section{Bayes' Rule for Classification}

\defn{Bayes Classifier}{
Given a new point $\vect{x}$, predict:
\[
\hat{y} = \arg\max_y P(Y=y \mid X=\vect{x}) = \arg\max_y \frac{P(X=\vect{x} \mid Y=y)\, P(Y=y)}{P(X=\vect{x})}.
\]
Since $P(X=\vect{x})$ is constant across classes:
\[
\hat{y} = \arg\max_y P(X=\vect{x} \mid Y=y)\, P(Y=y).
\]
$P(Y=y)$ is the \textbf{prior} (estimated from class frequencies in training data); $P(X \mid Y=y)$ is the \textbf{class-conditional likelihood}.
}

\fact{The Bayes classifier minimizes the probability of misclassification given perfect knowledge of $P(Y \mid X)$. The resulting error is the \textbf{Bayes error rate} $R^* = \mathbb{E}[1 - \max_y P(Y=y \mid X)]$, a fundamental lower bound on classification error.}

\section{The Na\"{i}ve Bayes Assumption}

The class-conditional $P(X \mid Y=y)$ is a $d$-dimensional distribution. Estimating it directly requires exponentially many parameters (e.g., a table with $2^d$ entries for binary features). The key simplification:

\defn{Na\"{i}ve Bayes (Conditional Independence) Assumption}{
Features are assumed conditionally independent given the class label:
\[
P(X_1, \ldots, X_d \mid Y=y) = \prod_{j=1}^d P(X_j \mid Y=y).
\]
This reduces estimation to $d$ separate 1-D problems.
}

The full prediction rule becomes:
\[
\hat{y} = \arg\max_y P(Y=y) \prod_{j=1}^d P(X_j = x_j \mid Y=y).
\]

\rmkb{The ``na\"{i}ve'' assumption is almost always violated in practice --- (``likes cats'' and ``likes dogs'' are surely correlated). Yet Na\"{i}ve Bayes remains surprisingly effective, especially for text classification, because the decision boundary is more robust than the probabilities themselves.}

\section{Parameter Estimation}

\textbf{Categorical features:} Estimate $P(X_j = v \mid Y = y) = \frac{\#\{i : x_{ij} = v, y_i = y\}}{\#\{i : y_i = y\}}$ by counting.

\textbf{Gaussian Na\"{i}ve Bayes:} Assume $P(X_j \mid Y=y) = \mathcal{N}(\mu_{jy}, \sigma_{jy}^2)$ for continuous features. Estimate $\mu_{jy}$ and $\sigma_{jy}^2$ by class-conditional sample mean and variance.

\section{Laplace Smoothing}

\defn{Laplace Smoothing}{
If a feature value never appears in class $y$, the estimate $P(X_j = v \mid Y=y) = 0$ zeros out the entire product. Laplace smoothing adds a pseudocount $\alpha > 0$:
\[
P(X_j = v \mid Y=y) = \frac{\#\{i : x_{ij}=v, y_i=y\} + \alpha}{\#\{i : y_i=y\} + \alpha\,|\mathcal{V}_j|},
\]
where $|\mathcal{V}_j|$ is the number of distinct values of feature $j$. For $\alpha=1$ this is called \textbf{add-one smoothing}.
}

In practice, computations are done in log-space to avoid numerical underflow from multiplying many probabilities:
\[
\hat{y} = \arg\max_y \left[\log P(Y=y) + \sum_{j=1}^d \log P(X_j = x_j \mid Y=y)\right].
\]

% ------------------------------------------------------------------------------
\chapter{Decision Trees and Ensembles}
% ------------------------------------------------------------------------------

\section{Decision Trees}

\defn{Decision Tree}{
A decision tree is a hierarchical model that recursively partitions the feature space using axis-aligned binary splits. Each internal node tests a feature $x_j$ against a threshold $t$; each leaf stores a predicted value (majority class or mean target).
}

A tree with $M$ leaves partitions $\R^d$ into $M$ rectangular regions $R_1, \ldots, R_M$. Predictions are constant within each region:
\[
f(\vect{x}) = \sum_{m=1}^M c_m\, \mathbf{1}[\vect{x} \in R_m].
\]

\section{Splitting Criteria}

Greedy tree learning chooses the split $(j, t)$ that maximizes the decrease in impurity.

\defn{Entropy and Information Gain}{
For a set $S$ with class proportions $p_1, \ldots, p_K$, the \textbf{Shannon entropy} is:
\[
H(S) = -\sum_{k=1}^K p_k \log_2 p_k.
\]
The \textbf{information gain} of splitting $S$ on feature $j$ at threshold $t$ is:
\[
\mathrm{IG}(S, j, t) = H(S) - \frac{|S_L|}{|S|} H(S_L) - \frac{|S_R|}{|S|} H(S_R),
\]
where $S_L, S_R$ are the left and right subsets after the split. The ID3/C4.5 algorithm greedily maximizes information gain.
}

\defn{Gini Impurity}{
\[
G(S) = 1 - \sum_{k=1}^K p_k^2 = \sum_{k \neq k'} p_k p_{k'}.
\]
$G(S) = 0$ for a pure node (one class); $G(S) = 1 - 1/K$ for maximum impurity. CART uses Gini impurity.
}

\prop{For regression trees, the splitting criterion minimizes the weighted sum of within-region MSE:
\[
\frac{1}{|S_L|}\sum_{\vect{x}_i \in S_L}(y_i - \bar{y}_{S_L})^2 + \frac{1}{|S_R|}\sum_{\vect{x}_i \in S_R}(y_i - \bar{y}_{S_R})^2.
\]}

\section{Overfitting and Pruning}

A fully-grown tree (one leaf per training point) memorizes training data. \textbf{Cost-complexity pruning} replaces subtrees with leaves while the gain is below a threshold. Alternatively, restrict tree depth, minimum leaf size, or minimum impurity decrease.

\section{Random Forests}

\defn{Bagging (Bootstrap Aggregating)}{
Train $B$ trees on independently bootstrapped datasets $\mathcal{D}_1^*, \ldots, \mathcal{D}_B^*$ (each drawn with replacement from $\mathcal{D}$). Average their predictions (regression) or take majority vote (classification):
\[
f_{\mathrm{bag}}(\vect{x}) = \frac{1}{B}\sum_{b=1}^B f_b(\vect{x}).
\]
Bagging reduces variance without increasing bias.
}

\defn{Random Forest}{
Extend bagging by also randomly subsampling $m \approx \sqrt{d}$ features at each split. This decorrelates the trees, further reducing variance. Random forests are robust, parallelizable, and provide out-of-bag error estimates.
}

\section{Boosting}

Boosting constructs an additive model $f(x) = \sum_{k=1}^K \alpha_k g(x; \gamma_k)$ by \textbf{sequentially} fitting weak learners $g$ (typically shallow trees), each one correcting the errors of its predecessors.

\defn{Forward Stagewise Additive Modeling}{
Initialize $f_0(x) = 0$. At each step $k = 1, \ldots, K$:
\begin{enumerate}
    \item Fit: $(\alpha_k, \gamma_k) = \arg\min_{\alpha, \gamma} \sum_{i=1}^n L\bigl(y_i,\, f_{k-1}(x_i) + \alpha g(x_i; \gamma)\bigr)$.
    \item Update: $f_k(x) = f_{k-1}(x) + \alpha_k g(x; \gamma_k)$.
\end{enumerate}
}

\defn{AdaBoost}{
For binary classification $y_i \in \{+1,-1\}$ with exponential loss $L(y, f) = \exp(-yf)$, the forward stagewise procedure reduces to AdaBoost:
\begin{enumerate}
    \item Initialize sample weights $u_i^{(1)} = 1$.
    \item For $k=1,\ldots,K$: train $g_k$ on reweighted data, compute weighted error $\varepsilon_k = \frac{\sum_i u_i^{(k)} \mathbf{1}[y_i \neq g_k(x_i)]}{\sum_i u_i^{(k)}}$, set $\alpha_k = \frac{1}{2}\log\frac{1-\varepsilon_k}{\varepsilon_k}$, then update $u_i^{(k+1)} = u_i^{(k)} \exp(-\alpha_k y_i g_k(x_i))$.
    \item Output $f(x) = \mathrm{sign}\!\left(\sum_k \alpha_k g_k(x)\right)$.
\end{enumerate}
Misclassified points receive higher weight in the next iteration; correctly classified points receive lower weight.
}

\prop{Under exponential loss, the stagewise procedure exactly corresponds to fitting each weak learner $g_k$ on the reweighted training set, with weights $u_i^{(k)} = \exp(-y_i f_{k-1}(x_i))$.}

\defn{Gradient Boosting}{
For a general differentiable loss $L$, at step $k$ compute the \textbf{pseudo-residuals}
\[
r_i = -\left.\frac{\pd{L(y_i, f(x_i))}{f(x_i)}}{}\right|_{f=f_{k-1}},
\]
and fit $g_k$ to the pseudo-residuals. For squared error loss, this is exactly fitting the residuals $r_i = y_i - f_{k-1}(x_i)$.
}

\rmkb{XGBoost and LightGBM are highly optimized gradient boosting implementations that dominate tabular data competitions. They add second-order Taylor expansions of the loss, regularization on leaf weights, and efficient histogram-based split finding.}

% ==============================================================================
% PART III: UNSUPERVISED LEARNING
% ==============================================================================

\part{Unsupervised Learning}

% ==============================================================================
\chapter{Clustering}
% ==============================================================================

\section{The Unsupervised Setting}

In supervised learning we are given labeled pairs $(x_i, y_i)$.  In the
\textit{unsupervised} setting we observe only $\{x_1,\dots,x_n\}$ and must
discover structure without any target signal.  Clustering is the task of
grouping examples so that points within the same cluster are similar and
points in different clusters are dissimilar.

\defn{Clustering}{
Given data $\{x_1,\dots,x_n\} \subset \R^d$ and a number of clusters $k$,
find a partition $\{S_1,\dots,S_k\}$ of the data and representative
centroids $\{c_1,\dots,c_k\}$ such that the \textbf{intra-cluster distortion}
is minimised:
\[
    \min_{\{S_j\},\,\{c_j\}}
    \sum_{j=1}^{k}\sum_{x\in S_j} \norm{x - c_j}^2.
\]
}

\section{K-Means Algorithm}

K-means (Steinhaus 1950 / Lloyd 1957) is an iterative algorithm that
alternates between two steps.

\fact{
\textbf{K-Means Algorithm.}  Input: $\{x_1,\dots,x_n\}$, $k$.\\[2pt]
\textbf{Initialise:} choose $c_1,\dots,c_k$ (e.g.\ random sample or
K-means++ — see below).\\[2pt]
\textbf{Repeat until centroids do not change:}
\begin{enumerate}
  \item \emph{Assignment step.} For every $x_i$ set
        $j^* = \arg\min_{j=1:k} \norm{x_i - c_j}^2$ and assign
        $S_{j^*} \leftarrow S_{j^*} \cup \{x_i\}$.
  \item \emph{Update step.} For $j = 1,\dots,k$ recompute
        $c_j \leftarrow \dfrac{1}{|S_j|}\displaystyle\sum_{x\in S_j} x$.
\end{enumerate}
}

An equivalent continuous objective reformulates the problem as:
\[
    \min_{c_1,\dots,c_k}\; f(c_1,\dots,c_k)
    = \frac{1}{2}\sum_{i=1}^{n} \min_{j=1:k} \norm{x_i - c_j}^2.
\]
The gradient update on centroid $c_j$ with learning rate $\eta$ gives the
SGD K-means rule:
\[
    c_j \;\leftarrow\; (1-\eta)\,c_j + \eta\,x_i
    \quad (\text{for }x_i\text{ whose closest centroid is }c_j).
\]

\section{Convergence}

\thm{K-Means Convergence}{
K-means always terminates.  At every iteration the objective
$\sum_j\sum_{x\in S_j}\norm{x-c_j}^2$ is non-increasing:
\begin{itemize}
  \item The assignment step only moves each $x$ to a closer (or equally
        close) centroid.
  \item The update step sets $c_j$ to the (unique) minimiser of
        $\sum_{x\in S_j}\norm{x-c_j}^2$, namely the mean.
\end{itemize}
Since there are finitely many possible partitions, the algorithm must stop.
}

K-means converges to a \emph{local} optimum; the global optimum is NP-hard
in general.  Running the algorithm multiple times with different
initialisations mitigates this.

\section{Initialisation: K-Means++}

\defn{K-Means++ Initialisation}{
Let $d_j(x)$ be the distance from $x$ to the closest of the $j$ centroids
already chosen.
\begin{enumerate}
  \item Choose $c_1$ uniformly at random from $\{x_1,\dots,x_n\}$.
  \item For $j = 2,\dots,k$: choose $c_j$ with probability proportional to
        $d_{j-1}(x)^2$, i.e.\
        \[
            P(x = c_j)
            = \frac{d_{j-1}(x)^2}{\sum_{i=1}^{n} d_{j-1}(x_i)^2}.
        \]
\end{enumerate}
This ensures each new centroid is likely to be far from all previous ones,
providing better coverage of the data.
}

\section{Connection to Gaussian Mixture Models}

Consider a uniform mixture of $k$ spherical Gaussians with means at the
K-means centroids:
\[
    P(x) \propto \frac{1}{k}\sum_{j=1}^{k} \exp\!\bigl(-\norm{x-c_j}^2\bigr).
\]
Because of exponential decay, only the closest centroid $z(x)
= \arg\min_{c_j}\norm{x-c_j}^2$ contributes appreciably:
\[
    P(X) \approx \mathrm{const}\times
    \exp\!\Bigl(-\sum_{i=1}^{n}\norm{x_i - z(x_i)}^2\Bigr).
\]
Hence maximising the likelihood under this model is equivalent to minimising
the K-means objective — K-means performs \emph{hard} EM on a Gaussian
mixture (vs.\ soft EM which maintains probabilistic assignments).

\section{Limitations of K-Means}

\rmkb{
\begin{itemize}
  \item K-means assumes \textbf{spherical clusters} of similar size and
        density.  It fails on elongated, non-convex, or heterogeneous
        clusters.
  \item The number of clusters $k$ must be specified in advance.  A common
        heuristic is the \emph{elbow method}: plot the objective vs.\ $k$
        and pick the ``knee'' of the curve.
  \item Because the partitioning space is discrete but large, a greedy
        algorithm may find a poor local minimum; running with multiple
        random restarts or K-means++ is recommended.
\end{itemize}
}

% ==============================================================================
\chapter{Dimensionality Reduction}
\label{ch:dimred}
% ==============================================================================

\section{Motivation}

High-dimensional data often lies near a lower-dimensional manifold.
Dimensionality reduction finds a mapping $x \in \R^d \to z \in \R^q$ with
$q \ll d$ that preserves as much information as possible.  Uses include
visualisation, removing redundant/correlated features, and reducing
computational cost.

\section{Principal Component Analysis}

\defn{PCA}{
\textbf{Principal Component Analysis (PCA)} finds the $q$ orthonormal
directions (principal components) in $\R^d$ that capture the most variance
in the data.  Equivalently, the projection onto these $q$ directions
minimises the mean reconstruction error.
}

The two interpretations are equivalent:

\begin{center}
\begin{tabular}{ll}
\toprule
\textbf{Maximum variance} & maximise variance of projections \\
\textbf{Minimum reconstruction error} & minimise $\sum_i \norm{x_i - \hat x_i}^2$ \\
\bottomrule
\end{tabular}
\end{center}

\subsection*{Procedure (Variance View)}

\begin{enumerate}
  \item \textbf{Centre the data:} $X \leftarrow X - \bar{X}$ where
        $\bar{X} = \frac{1}{n}\sum_{i=1}^n x_i$.
  \item Find the $q$ eigenvectors of $X^\top X$ with the $q$ largest
        eigenvalues; call them $v_1,\dots,v_q$.
  \item Project: $z_i = [v_1^\top x_i,\dots,v_q^\top x_i]^\top\in\R^q$.
\end{enumerate}

\subsection*{Derivation (First Principal Component)}

For a unit vector $v$ ($\norm{v}_2 = 1$) the variance of the projections
$Xv$ is (after centring, $\bar{Xv}=0$):
\[
    \mathrm{Var}(Xv)
    = \sum_i (v^\top x_i)^2
    = \norm{Xv}_2^2
    = v^\top X^\top X\, v.
\]
We want:
\[
    \max_{v:\,\norm{v}_2=1} v^\top X^\top X\, v.
\]
Introducing a Lagrange multiplier $\lambda$ and setting the gradient to zero:
\[
    \nabla\bigl[v^\top X^\top X v - \lambda\,v^\top v\bigr]
    = 2(X^\top X - \lambda I)v = 0
    \;\Longrightarrow\;
    \boxed{X^\top X\, v = \lambda v.}
\]
The maximum is attained at the eigenvector with the \emph{largest}
eigenvalue.  The $q$ principal components are the top-$q$ eigenvectors.

\prop{Variance Explained}{
If we keep the first $q$ principal components with eigenvalues
$\lambda_1 \ge \lambda_2 \ge \cdots$, the fraction of total variance
retained is
\[
    \frac{\sum_{k=1}^{q}\lambda_k}{\sum_{k=1}^{p}\lambda_k}.
\]
Plot this vs.\ $q$ (a \emph{scree plot}) to choose $q$.
}

\section{PCA via Singular Value Decomposition}

\defn{SVD}{
Every matrix $M \in \R^{n\times d}$ admits a factorisation
\[
    M = U\,\Sigma\,V^\top
\]
where $U\in\R^{n\times n}$ (columns = eigenvectors of $MM^\top$),
$\Sigma\in\R^{n\times d}$ (diagonal, singular values $\sigma_i = \sqrt{\lambda_i}$),
and $V\in\R^{d\times d}$ (columns = eigenvectors of $M^\top M$).
}

For PCA with centred data matrix $X$: compute $X = U\Sigma V^\top$; the
columns of $V$ are the principal directions, and $XV = U\Sigma$ are the
principal component scores.

\section{Autoencoders: Non-linear PCA}

A linear autoencoder with tied weights ($W$ and $W^\top$) trained to
minimise reconstruction error $\sum_{x\in D}\norm{x - \hat x}^2$ recovers
the PCA subspace.  A \textbf{deep autoencoder} replaces the linear layers
with non-linear networks (encoder $f_\theta$ and decoder $g_\phi$):

\[
    z = f_\theta(x) \in \R^k, \qquad
    \hat x = g_\phi(z) \in \R^d, \qquad
    \min_{\theta,\phi}\sum_{x\in D}\norm{x - g_\phi(f_\theta(x))}^2.
\]

The encoder maps $x\in\R^d$ to a low-dimensional code $z\in\R^k$ ($k\ll d$);
the decoder reconstructs $\hat x$.  Deep autoencoders can capture
non-linear manifold structure that PCA misses.

\rmkb{
\textbf{Autoencoder as pre-training.}  When labelled data is scarce, one
can (1) train an autoencoder on the large unlabelled corpus to learn a rich
latent representation, then (2) attach a classifier head to the encoder and
fine-tune on the small labelled set.  This is an early form of
\emph{self-supervised} pre-training.
}

\section{Limitations of PCA}

\begin{itemize}
  \item PCA is a \textbf{linear} projection; it fails when data lies on a
        non-linear manifold.
  \item It assumes approximately Gaussian-distributed features.
  \item Sensitive to scale; features should be standardised before applying
        PCA.
\end{itemize}

% ==============================================================================
\chapter{Learning Paradigms}
% ==============================================================================

\section{A Taxonomy of Learning}

Machine learning methods differ primarily in what supervision signal is
available during training.

\defn{Supervised Learning}{
Training data consists of labeled pairs $\{(x_i, y_i)\}_{i=1}^n$.
The goal is to learn a mapping $f: \mathcal{X} \to \mathcal{Y}$ that
generalises to unseen inputs.  Examples: linear/logistic regression, SVMs,
neural networks, decision trees.
}

\defn{Unsupervised Learning}{
Only unlabeled examples $\{x_i\}_{i=1}^n$ are available.  Goals include
discovering clusters, learning compact representations, estimating densities,
and generating new samples.  Examples: K-means, PCA, autoencoders, GMMs,
GANs.
}

\defn{Semi-Supervised Learning}{
A small labelled set $\mathcal{L} = \{(x_i, y_i)\}$ is supplemented by a
large unlabelled set $\mathcal{U} = \{x_j\}$.  The unlabelled data can
reveal the structure of $P(X)$, which (under the manifold or cluster
assumptions) helps improve decision boundaries learned from $\mathcal{L}$.
}

\defn{Self-Supervised Learning}{
Labels are derived automatically from the data itself.  A \emph{pretext
task} (e.g.\ predict a masked token, predict the next token, predict
image rotation) provides a supervised training signal without human
annotation.  The learned representations are then transferred to downstream
tasks.  Modern large language models (GPT, BERT) are trained this way.
}

\defn{Few-Shot Learning}{
At test time the model must generalise to new classes or tasks with only a
handful of labelled examples (the \emph{support set}).  Meta-learning
approaches (e.g.\ MAML, Prototypical Networks) aim to ``learn how to learn''
from prior tasks so adaptation to a new task requires minimal data.
}

\section{Transfer Learning}

In \textbf{transfer learning} a model pre-trained on a large source task is
adapted to a target task.

\begin{itemize}
  \item \textbf{Feature extraction:} freeze pre-trained weights; train only
        a new head on the target task.
  \item \textbf{Fine-tuning:} initialise from pre-trained weights; continue
        training all (or some) layers on the target task with a small
        learning rate.
\end{itemize}

Transfer learning is especially powerful when source and target distributions
overlap (e.g.\ ImageNet $\to$ medical imaging; GPT pre-training $\to$
task-specific RLHF fine-tuning).

\section{Reinforcement Learning}

In \textbf{reinforcement learning} (RL) an agent interacts with an
environment: at each step it observes a state $s_t$, takes an action $a_t$,
receives a reward $r_t$, and transitions to $s_{t+1}$.  The goal is to
learn a \emph{policy} $\pi(a|s)$ that maximises expected cumulative reward
$\mathbb{E}[\sum_t \gamma^t r_t]$, where $\gamma\in[0,1)$ is a discount
factor.  RL is covered in depth in Part~VIII.

% ==============================================================================
% PART IV: DEEP LEARNING
% ==============================================================================

\part{Deep Learning}

% ==============================================================================
\chapter{Neural Network Foundations}
% ==============================================================================

\section{The Single Neuron}

The basic computational unit is an artificial neuron.  Given input
$x \in \R^d$, it computes
\[
    f(x) = g(w^\top x + b),
\]
where $w \in \R^d$ is the weight vector, $b \in \R$ is the bias, and
$g:\R\to\R$ is a non-linear \emph{activation function}.  For the
squared-loss training objective $L = \tfrac{1}{2}(y - w^\top x - b)^2$,
the gradient with respect to $w$ is $\nabla_w L = -(y - w^\top x - b)\,x$.

\section{Multi-Layer Perceptrons}

\defn{Feedforward Neural Network}{
A feedforward network with $L$ layers is defined by the forward pass:
\[
    z^{[l]} = W^{[l]} a^{[l-1]} + b^{[l]}, \qquad
    a^{[l]} = g^{[l]}(z^{[l]}), \quad l = 1, \ldots, L,
\]
with $a^{[0]} = x$.  The network output is $\hat y = a^{[L]}$.
}

Stacking the layers, the full network computes
\[
    f(x) = g^{[L]}\!\bigl(W^{[L]}\,g^{[L-1]}\!\bigl(\cdots
    g^{[1]}(W^{[1]} x + b^{[1]})\cdots\bigr) + b^{[L]}\bigr).
\]

\prop{Necessity of Non-Linearity}{
If every $g^{[l]}$ is the identity, then
$f(x) = (W^{[L]}\cdots W^{[1]})x + b'$,
a single affine map.  Composing linear layers produces one linear layer;
depth buys expressiveness only with non-linearities.
}

\section{Activation Functions}

\begin{center}
\begin{tabular}{llll}
\hline
Name & Formula $g(z)$ & $g'(z)$ & Notes \\
\hline
Sigmoid     & $1/(1+e^{-z})$     & $g(1-g)$          & saturates; not zero-centred \\
Tanh        & $(e^z-e^{-z})/(e^z+e^{-z})$ & $1-g^2$ & zero-centred; still saturates \\
ReLU        & $\max(0,z)$        & $\mathbf{1}[z>0]$ & fast; \emph{dying neuron} risk \\
Leaky ReLU  & $\max(\varepsilon z,z)$ & $\varepsilon$ or $1$ & fixes dead neurons \\
\hline
\end{tabular}
\end{center}

\rmkb{
\textbf{Vanishing gradients.}  With sigmoid, $|\sigma'(z)|\le 0.25$,
so each layer multiplies the gradient by at most $0.25$.  For an $L$-layer
network the gradient magnitude decays as $(0.25)^L$.  ReLU avoids saturation
on the positive half-line and is the \textbf{default choice for hidden layers}.
}

\defn{Output Layer Design}{
Choose the final activation and loss to match the task:
\begin{itemize}
  \item Binary classification: sigmoid output,
        $L = -[y\log\hat y + (1-y)\log(1-\hat y)]$.
  \item Multi-class ($C$ classes): softmax $\hat y_k = e^{z_k}/\sum_j e^{z_j}$,
        cross-entropy loss.
  \item Regression: linear output (no activation), squared loss.
\end{itemize}
}

\fact{
The \emph{Universal Approximation Theorem} (Cybenko 1989) states that a
single hidden layer of sufficient width with a non-polynomial activation can
approximate any continuous function on a compact domain to arbitrary
accuracy.  Deeper networks achieve the same approximation with exponentially
fewer neurons.
}

% ==============================================================================
\chapter{Training Deep Networks}
% ==============================================================================

\section{Backpropagation}

Training minimises a loss $L$ by gradient descent.  Gradients are computed
via the chain rule in a single backward pass through the computation graph.

\defn{Backpropagation (Matrix Form)}{
Given the forward pass $z^{[l]} = W^{[l]} a^{[l-1]} + b^{[l]}$,
$a^{[l]} = g^{[l]}(z^{[l]})$, the backward pass from $l = L$ down to $1$:
\begin{align*}
    \nabla_{z^{[l]}} L &= \nabla_{a^{[l]}} L \odot (g^{[l]})'(z^{[l]}), \\
    \nabla_{W^{[l]}} L &= \nabla_{z^{[l]}} L \cdot (a^{[l-1]})^\top, \\
    \nabla_{a^{[l-1]}} L &= (W^{[l]})^\top \nabla_{z^{[l]}} L.
\end{align*}
Modern frameworks (PyTorch, JAX) implement this via \emph{automatic
differentiation}: they record a computation graph during the forward pass
and apply the chain rule in reverse automatically.
}

\section{Weight Initialisation}

\defn{Initialisation Schemes}{
\begin{itemize}
  \item \textbf{Xavier (Glorot)}: $W_{ij} \sim \mathcal{U}\!\left(
        -\sqrt{\tfrac{6}{n_\text{in}+n_\text{out}}},\;
        \sqrt{\tfrac{6}{n_\text{in}+n_\text{out}}}\right)$.  Designed for
        tanh/sigmoid to preserve gradient variance.
  \item \textbf{He}: $W_{ij} \sim \mathcal{N}(0,\,2/n_\text{in})$.
        Designed for ReLU.
\end{itemize}
}

\section{Batch Normalisation}

\defn{Batch Normalisation (Ioffe \& Szegedy 2015)}{
For mini-batch $\{x^{(i)}\}_{i=1}^m$ and feature $k$, let
$\mu_k = \tfrac{1}{m}\sum_i x_k^{(i)}$,
$\sigma_k^2 = \tfrac{1}{m}\sum_i (x_k^{(i)}-\mu_k)^2$.  Batch norm:
\[
    \hat x_k^{(i)} = \frac{x_k^{(i)}-\mu_k}{\sqrt{\sigma_k^2+\varepsilon}},
    \qquad y_k^{(i)} = \gamma_k \hat x_k^{(i)} + \beta_k,
\]
where $\gamma_k, \beta_k$ are learned per-feature scale and shift.
}

Batch norm keeps activation distributions stable across layers, allowing
higher learning rates and faster convergence.  It also acts as a regulariser.

\section{Dropout}

\defn{Dropout (Srivastava et al.\ 2014)}{
During training, independently zero each neuron with probability $p$;
scale surviving activations by $\tfrac{1}{1-p}$ (\emph{inverted dropout}).
At inference, use all neurons without scaling.
}

Dropout prevents co-adaptation of neurons and acts as implicit ensemble
averaging, reducing overfitting.  Typical values: $p = 0.2$--$0.5$.

\section{Optimisation}

\textbf{SGD with Momentum}:
\[
    v \leftarrow \beta v - \eta \nabla_\theta L, \quad
    \theta \leftarrow \theta + v.
\]

\defn{Adam (Kingma \& Ba 2015)}{
Maintains exponential moving averages of gradients $m_t$ and squared
gradients $v_t$:
\[
    m_t = \beta_1 m_{t-1} + (1-\beta_1) g_t, \quad
    v_t = \beta_2 v_{t-1} + (1-\beta_2) g_t^2.
\]
Bias-corrected estimates $\hat m_t = m_t/(1-\beta_1^t)$,
$\hat v_t = v_t/(1-\beta_2^t)$.  Update:
\[
    \theta_t = \theta_{t-1} - \eta\,\frac{\hat m_t}{\sqrt{\hat v_t}+\varepsilon}.
\]
Defaults: $\beta_1=0.9$, $\beta_2=0.999$, $\varepsilon=10^{-8}$.
}

\section{Learning Rate Schedules}

The non-convex DNN loss surface benefits from a large $\eta$ early
(escape saddle points) and small $\eta$ later (fine convergence).

\begin{itemize}
  \item \textbf{Step decay}: multiply $\eta$ by $\gamma < 1$ every $k$ epochs.
  \item \textbf{Cosine annealing}:
        $\eta_t = \eta_{\min} + \tfrac{1}{2}(\eta_{\max}-\eta_{\min})
        (1+\cos(\pi t/T))$.
  \item \textbf{Warmup}: linearly increase $\eta$ for the first few steps to
        stabilise early training (used in Transformers).
\end{itemize}

\section{Data Augmentation and Mixup}

\textbf{Data augmentation} enlarges the effective training set by applying
label-preserving transformations (random crop, horizontal flip, colour
jitter, rotation for images).

\textbf{Mixup} (Zhang et al.\ 2018) interpolates both inputs and soft labels:
\[
    \hat x = \lambda x_i + (1-\lambda) x_j, \quad
    \hat y = \lambda y_i + (1-\lambda) y_j, \quad
    \lambda \sim \mathrm{Beta}(\alpha,\alpha).
\]
Mixup improves generalisation and reduces overconfident predictions.

% ==============================================================================
\chapter{Convolutional Neural Networks}
% ==============================================================================

\section{The Convolution Operation}

\defn{2-D Convolution Layer}{
Given input volume $H \times W \times C_\text{in}$, a filter bank of $K$
filters each of size $F \times F \times C_\text{in}$ produces an output
volume of size $H' \times W' \times K$:
\[
    H' = \frac{H - F + 2P}{S} + 1, \qquad W' = \frac{W - F + 2P}{S} + 1,
\]
with padding $P$ and stride $S$.
}

Key properties:
\begin{itemize}
  \item \textbf{Weight sharing}: the same filter applies at every spatial
        location — a strong inductive bias for translation equivariance.
  \item \textbf{Local connectivity}: each output neuron depends only on an
        $F \times F$ region, giving sparse connections vs.\ fully-connected.
  \item \textbf{Depth stacking}: $n$ conv layers of size $F$ give effective
        receptive field of $n(F-1)+1$.
\end{itemize}

\section{Pooling}

\defn{Max Pooling}{
Partitions the spatial dimensions into non-overlapping $F \times F$ cells
and outputs the maximum in each cell.  Reduces spatial resolution by $F$,
providing translational invariance and reducing computation.
}

\section{Typical CNN Architecture}

\[
    [\mathrm{Conv} \to \mathrm{ReLU} \to \mathrm{Pool}]^{\times N}
    \to \mathrm{Flatten}
    \to [\mathrm{FC} \to \mathrm{ReLU}]^{\times M}
    \to \mathrm{FC} \to \mathrm{Softmax}.
\]

Early layers detect low-level features (edges, colours); deeper layers
compose these into high-level semantic concepts.

\section{Classic Architectures}

\textbf{LeNet-5} (LeCun 1998): earliest successful deep CNN for digit
recognition.  Architecture on 32$\times$32 input:
$\mathrm{C1}(6@28^2) \to \mathrm{S2}(6@14^2) \to \mathrm{C3}(16@10^2)
\to \mathrm{S4}(16@5^2) \to \mathrm{C5}(120) \to \mathrm{F6}(84)
\to \mathrm{Output}(10)$.

\textbf{VGGNet} (Simonyan \& Zisserman 2014): stacks of $3\times3$
convolutions and $2\times2$ max-pooling.  Three $3\times3$ layers give the
same $7\times7$ receptive field as one $7\times7$ filter but with fewer
parameters and more non-linearities.  VGG-16 has 138M parameters.

\textbf{GoogLeNet / InceptionNet} (Szegedy et al.\ 2014): \emph{Inception
modules} apply $1\times1$, $3\times3$, and $5\times5$ convolutions and
$3\times3$ max-pooling in parallel and concatenate along the channel
dimension.  $1\times1$ convolutions reduce channel dimensions cheaply
(\emph{bottleneck}).  Only 5M parameters.

\textbf{ResNet} (He et al.\ 2016): introduces \emph{residual (skip)
connections} to solve the degradation problem — deeper plain networks
having higher training error than shallower ones.

\defn{Residual Block}{
\[
    H(x) = F(x) + x,
\]
where $F(x)$ is the residual to be learned.  If the identity mapping is
optimal, $F$ is driven to zero.  The skip connection provides a gradient
highway to earlier layers, enabling networks of 50--152+ layers.
}

\prop{Bottleneck Residual Block}{
For ResNet-50/101/152 the block uses: $1\times1$ conv (reduce channels)
$\to$ $3\times3$ conv $\to$ $1\times1$ conv (restore channels), achieving
similar complexity to a $3\times3$ two-layer block but at greater depth.
}

% ==============================================================================
\chapter{Object Detection and Segmentation}
% ==============================================================================

\section{From Classification to Localisation}

Image classification assigns one label to the whole image.
\textbf{Object detection} simultaneously localises all objects with
\emph{bounding boxes} and classifies each one.

\defn{Bounding Box and IoU}{
A bounding box is parametrised as $(x_c, y_c, w, h)$.  Detection quality is
measured by Intersection over Union:
\[
    \mathrm{IoU}(\hat B, B^*) = \frac{|\hat B \cap B^*|}{|\hat B \cup B^*|}.
\]
A prediction is a true positive when $\mathrm{IoU} \geq 0.5$ (PASCAL VOC)
or averaged over thresholds 0.5--0.95 (COCO mAP).
}

\section{Two-Stage Detectors: Faster R-CNN}

A \textbf{Region Proposal Network (RPN)} slides a small convolutional head
over the backbone feature map, predicting objectness scores and bounding-box
offsets for a set of \emph{anchors} (fixed-scale, fixed-ratio boxes).
Proposals with high objectness are passed to an ROI-pooling layer that
extracts fixed-size feature vectors, which are then classified and refined
by detection heads.  Faster R-CNN achieves high accuracy but requires two
forward passes (RPN + detection).

\section{One-Stage Detectors: YOLO}

\textbf{YOLO} (Redmon et al.\ 2016) divides the image into an $S \times S$
grid; each cell predicts $B$ bounding boxes with confidence scores and $C$
class probabilities in a single forward pass, enabling real-time detection.

\section{Semantic Segmentation}

\defn{Encoder--Decoder Segmentation}{
Full per-pixel classification.  An encoder (convolutional backbone)
downsamples the feature map; a decoder upsamples via transposed convolutions
to the original resolution.  \emph{Skip connections} from encoder to decoder
preserve fine spatial detail.  U-Net is the canonical example.
}

% ==============================================================================
\chapter{Recurrent Neural Networks}
% ==============================================================================

\section{Sequence Modelling}

Many tasks involve variable-length sequential inputs and outputs.  Standard
feedforward networks require fixed-size inputs; RNNs process sequences by
maintaining a \emph{hidden state} that summarises the history.

\defn{Sequence Topologies}{
\begin{itemize}
  \item \textbf{Many-to-one}: whole sequence $\to$ single output
        (e.g.\ sentiment classification).
  \item \textbf{One-to-many}: single input $\to$ sequence (e.g.\ image captioning).
  \item \textbf{Many-to-many (aligned)}: per-step prediction (e.g.\ video labelling).
  \item \textbf{Seq2seq}: many-to-one encoder compresses input; one-to-many
        decoder generates output (e.g.\ machine translation).
\end{itemize}
}

\section{Vanilla RNN}

\defn{Elman RNN}{
At each time step $t$, with input $x_t \in \R^{d}$ and previous hidden state
$h_{t-1} \in \R^{h}$:
\begin{align*}
    h_t &= \tanh(W_{hh}\,h_{t-1} + W_{xh}\,x_t + b_h), \\
    y_t &= W_{hy}\,h_t + b_y.
\end{align*}
The \textbf{same} weight matrices $W_{hh}, W_{xh}, W_{hy}$ are reused at
every step (\emph{weight sharing across time}).
}

\section{Backpropagation Through Time}

\defn{BPTT}{
The gradient of the total loss $J_T = \sum_{t=1}^T L_t$ with respect to
shared weights is summed over time steps:
\[
    \frac{\partial J_T}{\partial W}
    = \sum_{i=1}^{T} \left.\frac{\partial J_T}{\partial W}\right|_{(i)}.
\]
The gradient at step $T$ with respect to an earlier hidden state $h_t$
requires the product
\[
    \frac{\partial J_T}{\partial h_t}
    = \frac{\partial J_T}{\partial h_T}\,W_{hh}^{T-t}
      \prod_{t < i \leq T}
      \mathrm{diag}\bigl[\tanh'(W_{hh} h_{i-1} + W_{xh} x_i)\bigr].
\]
}

\prop{Vanishing and Exploding Gradients}{
The factor $W_{hh}^{T-t}$ causes the gradient to:
\begin{itemize}
  \item \textbf{Vanish}: when singular values of $W_{hh}$ are $< 1$ —
        early time steps receive negligible gradient; no long-range learning.
  \item \textbf{Explode}: when singular values of $W_{hh}$ are $> 1$ —
        training diverges.
\end{itemize}
Fix for exploding gradients: \textbf{gradient clipping} — if
$\norm{g} \geq \tau$, set $g \leftarrow (\tau/\norm{g})\,g$.  Fix for
vanishing gradients: \textbf{gated architectures} (LSTM, GRU).
}

\section{Long Short-Term Memory}

\defn{LSTM (Hochreiter \& Schmidhuber 1997)}{
An LSTM maintains a \emph{cell state} $C_t$ (long-term memory) and a
\emph{hidden state} $H_t$ (short-term memory).  At each step, given the
concatenation $[H_{t-1}; X_t]$:
\begin{align*}
    F_t &= \sigma\!\bigl(W_F [H_{t-1}; X_t] + b_F\bigr) & &\text{(forget gate)} \\
    I_t &= \sigma\!\bigl(W_I [H_{t-1}; X_t] + b_I\bigr) & &\text{(input gate)} \\
    \tilde C_t &= \tanh\!\bigl(W_C [H_{t-1}; X_t] + b_C\bigr) & &\text{(candidate cell)} \\
    C_t &= F_t \odot C_{t-1} + I_t \odot \tilde C_t & &\text{(cell update)} \\
    O_t &= \sigma\!\bigl(W_O [H_{t-1}; X_t] + b_O\bigr) & &\text{(output gate)} \\
    H_t &= O_t \odot \tanh(C_t) & &\text{(hidden state)}
\end{align*}
}

The cell state flows through time with only element-wise operations
(forget + add), providing a gradient highway and enabling learning of
dependencies across hundreds of steps.

\rmkb{
\textbf{GRU} (Cho et al.\ 2014) merges cell and hidden state into one vector
and uses two gates (update gate $z_t$, reset gate $r_t$), matching LSTM
performance with fewer parameters.
}

\section{Bidirectional and Multi-Layer RNNs}

\defn{Bidirectional RNN}{
Two RNNs process the sequence in forward and backward directions; hidden
states are concatenated at each step:
\[
    \vec h_t = \mathrm{RNN}_{\mathrm{FW}}(\vec h_{t-1},\, x_t), \quad
    \overleftarrow{h}_t = \mathrm{RNN}_{\mathrm{BW}}(\overleftarrow{h}_{t+1},\, x_t), \quad
    h_t = [\vec h_t;\, \overleftarrow{h}_t].
\]
Each position's representation incorporates both past and future context.
Transformers largely supersede bidirectional RNNs due to parallelism.
}

\textbf{Multi-layer RNNs} feed the hidden states of layer $\ell$ as inputs
to layer $\ell+1$, enabling hierarchical temporal feature extraction.

% ==============================================================================
\chapter{Attention and Transformers}
% ==============================================================================

\section{Scaled Dot-Product Attention}

\defn{Attention}{
Given a query $q \in \R^{d_k}$, keys $K \in \R^{n \times d_k}$, and values
$V \in \R^{n \times d_v}$:
\[
    \mathrm{Attention}(q, K, V)
    = \mathrm{softmax}\!\left(\frac{q K^\top}{\sqrt{d_k}}\right) V.
\]
The factor $1/\sqrt{d_k}$ prevents dot products from saturating the softmax.
}

In \textbf{self-attention}, queries, keys, and values are all derived from
the same input sequence $X$ via learned projections:
\[
    Q = X W_Q, \quad K = X W_K, \quad V = X W_V.
\]
Every token attends to every other token simultaneously in $O(1)$ layers,
unlike RNNs which require $O(T)$ sequential steps.

\section{Multi-Head Attention}

\defn{Multi-Head Attention}{
Run $h$ independent attention heads in parallel, each with projections
$W_Q^i, W_K^i, W_V^i \in \R^{d_\text{model} \times d_k}$, then concatenate
and project:
\[
    \mathrm{MultiHead}(Q,K,V)
    = \mathrm{Concat}(\mathrm{head}_1, \ldots, \mathrm{head}_h)\,W_O,
    \quad
    \mathrm{head}_i = \mathrm{Attention}(QW_Q^i,\,KW_K^i,\,VW_V^i).
\]
Different heads can capture syntactic, semantic, and positional relationships
simultaneously.
}

\section{The Transformer Architecture}

\defn{Transformer (Vaswani et al.\ 2017)}{
An encoder--decoder architecture without recurrence.

\textbf{Encoder block} (repeated $N$ times):
\begin{enumerate}
  \item Multi-head self-attention $\to$ Add \& LayerNorm.
  \item Position-wise FFN $\to$ Add \& LayerNorm.
\end{enumerate}

\textbf{Decoder block} (repeated $N$ times):
\begin{enumerate}
  \item Masked multi-head self-attention (causal; prevents attending to future)
        $\to$ Add \& LayerNorm.
  \item Multi-head cross-attention (queries from decoder; keys, values from
        encoder) $\to$ Add \& LayerNorm.
  \item Position-wise FFN $\to$ Add \& LayerNorm.
\end{enumerate}
}

\textbf{Feed-Forward sublayer}:
\[
    \mathrm{FFN}(x) = \max(0,\, x W_1 + b_1)\,W_2 + b_2.
\]

\textbf{Layer Normalisation}: for each token $h_k$ normalise over the feature
dimension, then apply learned scale $\gamma$ and bias $\beta$:
\[
    \mathrm{LN}(h_k)
    = \frac{h_k - \mu_k}{\sqrt{\sigma_k + \varepsilon}} \cdot \gamma + \beta.
\]

\defn{Positional Encoding}{
Since Transformers have no recurrence, token order is injected by adding
sinusoidal positional embeddings to the input:
\[
    \mathrm{PE}_{(pos,\,2i)}   = \sin\!\left(\frac{pos}{10000^{2i/d}}\right),
    \quad
    \mathrm{PE}_{(pos,\,2i+1)} = \cos\!\left(\frac{pos}{10000^{2i/d}}\right).
\]
Alternatively, learned positional embeddings are used in BERT/GPT.
}

\section{Pre-Trained Transformers}

\textbf{BERT} (Devlin et al.\ 2019) is a bidirectional Transformer encoder
pre-trained with two self-supervised tasks:
\begin{itemize}
  \item \textbf{Masked Language Modelling (MLM)}: randomly mask 15\% of
        tokens; predict the masked tokens.
  \item \textbf{Next Sentence Prediction (NSP)}: binary prediction of whether
        two sentences are consecutive.
\end{itemize}
BERT-base: 12 layers, hidden size 768, 12 attention heads, 110M parameters.
BERT-large: 24 layers, hidden size 1024, 16 heads, 340M parameters.
Training corpus: Wikipedia (2.5B words) + BookCorpus (0.8B words).
Fine-tuning: attach a task head to the pre-trained encoder; update all
parameters on labelled downstream data (\emph{pretrain once, fine-tune many
times}).

\textbf{GPT} (Radford et al.\ 2018) uses a causal (left-to-right) Transformer
\emph{decoder} pre-trained with standard language modelling:
\[
    L = \frac{1}{T}\sum_{t=1}^{T} L_{\mathrm{CE}}(y_t,\,\hat y_t).
\]
GPT-3 (Brown et al.\ 2020): 175B parameters, trained on 300B tokens.
It exhibits \textbf{in-context learning}: given task demonstrations in the
prompt, it generalises with no gradient updates (a purely forward-pass
phenomenon).

\rmkb{
\textbf{Transformer family.}
Encoder-only (BERT, RoBERTa): bidirectional context; excel at classification
and NLU. \\
Decoder-only (GPT, PaLM): auto-regressive; natural language generation. \\
Encoder-decoder (T5, BART): seq2seq tasks (translation, summarisation).
}

% ==============================================================================
% PART V: GENERATIVE MODELS
% ==============================================================================

\part{Generative Models}

% ==============================================================================
\chapter{Variational Autoencoders}
% ==============================================================================

\section{The Generative Modelling Goal}

A \textbf{generative model} learns the data distribution $p(x)$ so that one
can (i) evaluate likelihood and (ii) draw new samples $x \sim p(x)$.
Standard autoencoders compress data but impose no structure on the latent
space, making the generative direction unreliable.

\section{Evidence Lower Bound}

\defn{VAE (Kingma \& Welling 2014)}{
Assume a latent variable model $p_\theta(x,z)=p_\theta(x|z)\,p(z)$ with
Gaussian prior $p(z)=\mathcal{N}(0,I)$.  The marginal
$p_\theta(x)=\int p_\theta(x|z)\,p(z)\,dz$ is intractable.  Introducing an
approximate posterior $q_\phi(z|x)$, the \emph{Evidence Lower Bound} is:
\[
    \log p_\theta(x) \;\geq\;
    \mathcal{L}(\theta,\phi;\,x)
    = \underbrace{\mathbb{E}_{q_\phi(z|x)}[\log p_\theta(x|z)]}_{\text{reconstruction}}
    \;-\; \underbrace{D_{\mathrm{KL}}(q_\phi(z|x)\,\|\,p(z))}_{\text{regularisation}}.
\]
}

\prop{ELBO Derivation}{
By Jensen's inequality applied to the concave $\log$:
\[
    \log p_\theta(x)
    = \log \mathbb{E}_{q_\phi}\!\left[\frac{p_\theta(x|z)\,p(z)}{q_\phi(z|x)}\right]
    \geq \mathbb{E}_{q_\phi}\!\left[\log \frac{p_\theta(x|z)\,p(z)}{q_\phi(z|x)}\right]
    = \mathcal{L}.
\]
The gap is $D_{\mathrm{KL}}(q_\phi(z|x)\|p_\theta(z|x)) \geq 0$.
}

\section{The Reparametrisation Trick}

Sampling $z \sim q_\phi(z|x) = \mathcal{N}(\mu_\phi(x), \sigma_\phi^2(x)I)$
is non-differentiable.  Reparametrise:
\[
    z = \mu_\phi(x) + \sigma_\phi(x) \odot \varepsilon,
    \quad \varepsilon \sim \mathcal{N}(0,I).
\]
Randomness is now in $\varepsilon$ (no parameters), so gradients flow through
$\mu_\phi$ and $\sigma_\phi$ via standard backpropagation.

\rmkb{
\textbf{Closed-form KL} (diagonal Gaussian):
$D_{\mathrm{KL}}(\mathcal{N}(\mu,\sigma^2)\|\mathcal{N}(0,1))
= -\tfrac{1}{2}\sum_j(1 + \log\sigma_j^2 - \mu_j^2 - \sigma_j^2)$.
}

% ==============================================================================
\chapter{Generative Adversarial Networks}
% ==============================================================================

\section{The GAN Minimax Game}

\defn{GAN (Goodfellow et al.\ 2014)}{
Train two networks adversarially:
\begin{itemize}
  \item \textbf{Generator} $G_\theta:\mathcal{Z}\to\mathcal{X}$: maps noise
        $z\sim p(z)$ to fake samples.
  \item \textbf{Discriminator} $D_\phi:\mathcal{X}\to[0,1]$: classifies real
        vs.\ generated.
\end{itemize}
Minimax objective:
\[
    \min_\theta\max_\phi\;
    \mathbb{E}_{x\sim p_\text{data}}[\log D_\phi(x)]
    + \mathbb{E}_{z\sim p(z)}[\log(1-D_\phi(G_\theta(z)))].
\]
}

\prop{Optimal Discriminator and Equilibrium}{
At fixed $G$, the optimal discriminator is
$D^*(x) = p_\text{data}(x)/(p_\text{data}(x)+p_G(x))$.
Substituting, the minimax problem reduces to minimising the
Jensen--Shannon divergence:
$\min_\theta 2\,D_{\mathrm{JS}}(p_\text{data}\|p_G)-\log 4$.
Global minimum $p_G = p_\text{data}$.
}

\section{Practical Variants}

Training alternates: update $D$ to distinguish real from fake, then update
$G$ to fool $D$ (using the non-saturating loss $-\mathbb{E}[\log D_\phi(G_\theta(z))]$
to avoid early vanishing gradients).

\textbf{WGAN} (Arjovsky et al.\ 2017) uses the Wasserstein-1 distance and a
1-Lipschitz critic.  With gradient penalty:
\[
    L_{\mathrm{GP}} = \mathbb{E}[D(\tilde x)] - \mathbb{E}[D(x)]
    + \lambda\,\mathbb{E}[(\norm{\nabla_{\hat x} D(\hat x)}_2-1)^2].
\]

\textbf{Conditional GAN}: condition both $G$ and $D$ on auxiliary label $y$
to control the class or attributes of generated samples.

% ==============================================================================
\chapter{Diffusion Models}
% ==============================================================================

\section{Forward Diffusion Process}

\defn{DDPM (Ho et al.\ 2020)}{
A fixed \emph{forward process} gradually adds Gaussian noise over $T$ steps:
\[
    q(x_t \mid x_{t-1})
    = \mathcal{N}\!\bigl(\sqrt{1-\beta_t}\,x_{t-1},\;\beta_t I\bigr).
\]
Using $\bar\alpha_t = \prod_{s=1}^t(1-\beta_s)$, the marginal is:
\[
    q(x_t \mid x_0) = \mathcal{N}\!\bigl(\sqrt{\bar\alpha_t}\,x_0,\;(1-\bar\alpha_t)I\bigr),
\]
so $x_t = \sqrt{\bar\alpha_t}\,x_0 + \sqrt{1-\bar\alpha_t}\,\varepsilon$,
$\varepsilon\sim\mathcal{N}(0,I)$.
}

\section{Reverse Process and Training}

A neural network $\varepsilon_\theta(x_t,t)$ predicts the noise.  The
simplified training loss is:
\[
    \mathcal{L}_\text{simple}
    = \mathbb{E}_{t,\,x_0,\,\varepsilon}\!\left[
      \norm{\varepsilon
      - \varepsilon_\theta\!\bigl(\sqrt{\bar\alpha_t}\,x_0
        + \sqrt{1-\bar\alpha_t}\,\varepsilon,\;t\bigr)}_2^2\right].
\]

To generate: start from $x_T\sim\mathcal{N}(0,I)$ and iterate
\[
    x_{t-1}
    = \frac{1}{\sqrt{\alpha_t}}\!\left(x_t
      - \frac{1-\alpha_t}{\sqrt{1-\bar\alpha_t}}\varepsilon_\theta(x_t,t)\right)
    + \sigma_t z, \quad z\sim\mathcal{N}(0,I).
\]

\rmkb{
\textbf{Classifier-Free Guidance}: train $\varepsilon_\theta(x_t,t,c)$
jointly with the unconditional model (randomly drop $c$ during training).
At inference, amplify the conditioning signal:
$\tilde\varepsilon = \varepsilon_\theta(x_t,t) + w(\varepsilon_\theta(x_t,t,c)-\varepsilon_\theta(x_t,t))$
with guidance scale $w>1$.  Higher $w$ trades sample diversity for
prompt fidelity (used in Stable Diffusion, DALL-E 2).
}

% ==============================================================================
% PART VI: LARGE LANGUAGE MODELS
% ==============================================================================

\part{Large Language Models}

% ==============================================================================
\chapter{Language Models and Pretraining}
% ==============================================================================

\section{Language Modelling}

\defn{Language Model}{
A language model assigns probabilities to token sequences via the chain rule:
\[
    P(x_1,\ldots,x_T) = \prod_{t=1}^T P(x_t \mid x_1,\ldots,x_{t-1}).
\]
Training minimises cross-entropy (equivalently, maximises log-likelihood):
\[
    L = -\frac{1}{|\mathcal{D}|}\sum_t \log P_\theta(x_t \mid x_{<t}).
\]
\emph{Perplexity}: $\mathrm{PPL} = \exp(L)$.  Lower is better.
}

\section{Subword Tokenisation}

Raw text is mapped to integers via a subword tokeniser.
\textbf{Byte-Pair Encoding (BPE)}: iteratively merge the most frequent
adjacent byte/character pair, building a vocabulary of $|V|$ subword units
(typically 30k--100k).

\section{Transformer LM Architectures}

\begin{itemize}
  \item \textbf{Causal (decoder-only)} — GPT family: masked self-attention;
        trained on next-token prediction; natural text generation.
  \item \textbf{Masked (encoder-only)} — BERT family: bidirectional attention;
        trained on masked token prediction; strong NLU representations.
  \item \textbf{Encoder-decoder} — T5, BART: bidirectional encoder plus
        causal decoder; suited to seq2seq tasks.
\end{itemize}

\section{Scaling Laws}

\prop{Chinchilla Scaling Law (Hoffmann et al.\ 2022)}{
Loss improves as a power law in model size $N$ and training tokens $D$:
\[
    L(N,D) \approx \frac{A}{N^\alpha} + \frac{B}{D^\beta} + L_\infty,
    \quad \alpha\approx\beta\approx 0.5.
\]
For compute budget $C\approx 6ND$, the optimal allocation satisfies
$N^*\approx D^*\propto\sqrt{C}$: a smaller model trained on more tokens
outperforms a larger model trained on fewer.
}

% ==============================================================================
\chapter{Instruction Finetuning and Alignment}
% ==============================================================================

\section{Supervised Fine-Tuning}

Pre-trained LLMs predict the next token but do not inherently follow
instructions.  \textbf{Instruction fine-tuning} trains on curated
(prompt, response) pairs with standard cross-entropy, teaching the model
to be helpful, harmless, and honest.

\section{Reinforcement Learning from Human Feedback}

\defn{RLHF (Christiano et al.\ 2017; Ouyang et al.\ 2022)}{
\begin{enumerate}
  \item \textbf{Reward model}: rank model outputs; train $r_\phi(x,y)$ via
        \[
            L_{\mathrm{RM}} = -\mathbb{E}_{(x,\,y_w \succ y_l)}
            [\log\sigma(r_\phi(x,y_w)-r_\phi(x,y_l))].
        \]
  \item \textbf{RL fine-tuning} with PPO: maximise expected reward subject to
        a KL penalty that keeps the policy close to the SFT reference:
        \[
            \max_\theta\;\mathbb{E}_{x,\,y\sim\pi_\theta}[r_\phi(x,y)]
            - \beta\,D_{\mathrm{KL}}(\pi_\theta \| \pi_\mathrm{ref}).
        \]
\end{enumerate}
}

\section{Direct Preference Optimisation}

\defn{DPO (Rafailov et al.\ 2023)}{
Bypasses explicit reward model training.  The DPO loss:
\[
    L_{\mathrm{DPO}} = -\mathbb{E}_{(x,\,y_w,\,y_l)}\!\left[
      \log\sigma\!\left(
        \beta\log\frac{\pi_\theta(y_w|x)}{\pi_\mathrm{ref}(y_w|x)}
        - \beta\log\frac{\pi_\theta(y_l|x)}{\pi_\mathrm{ref}(y_l|x)}
      \right)
    \right].
\]
DPO is simpler, more stable, and often matches or exceeds RLHF quality.
}

% ==============================================================================
\chapter{In-Context Learning and Prompting}
% ==============================================================================

\section{In-Context Learning}

\defn{In-Context Learning (Brown et al.\ 2020)}{
Feed the model a prompt containing $k$ demonstrations
$(x_1,y_1),\ldots,(x_k,y_k)$ then a query $x_{k+1}$; the model predicts
$y_{k+1}$ via a \textbf{single forward pass} — no gradient updates.
Zero-shot: no demonstrations.  Few-shot: $k \geq 1$.
}

\section{Chain-of-Thought Prompting}

\textbf{Chain-of-Thought (CoT)} (Wei et al.\ 2022): include step-by-step
reasoning in few-shot exemplars.  Dramatically improves performance on
arithmetic, commonsense, and symbolic reasoning.

\textbf{Zero-Shot CoT}: append ``Let's think step by step.'' to the prompt;
elicits reasoning without examples.

\textbf{Self-Consistency}: sample $k$ reasoning chains independently;
take the majority-vote final answer.  Reduces variance and improves accuracy.

\rmkb{
\textbf{Least-to-Most Prompting}: decompose a complex problem into easier
subproblems; solve each sequentially, passing prior solutions into the
next context.  Enables compositional generalisation beyond training length.
}

% ==============================================================================
\chapter{Decoding Strategies}
% ==============================================================================

\section{Deterministic Decoding}

\defn{Greedy Decoding}{
$x_t = \arg\max_v P_\theta(v\mid x_{<t})$.  Fast but locally optimal;
can produce repetitive or suboptimal sequences.
}

\defn{Beam Search}{
Maintain $B$ (\emph{beam width}) partial hypotheses; expand each by all
vocabulary tokens; keep top-$B$ by cumulative log-probability.  Better for
structured tasks (translation, summarisation); tends to produce generic text
for open-ended generation.
}

\section{Stochastic Decoding}

\defn{Temperature Sampling}{
Sample from $P_T(v) \propto P(v)^{1/T}$.  $T<1$: peakier (more greedy);
$T>1$: flatter (more random).
}

\defn{Top-$k$ Sampling}{
Restrict to the $k$ highest-probability tokens, renormalise, sample.
}

\defn{Nucleus (Top-$p$) Sampling (Holtzman et al.\ 2020)}{
Restrict to the smallest vocabulary subset whose cumulative probability
$\geq p$ (e.g.\ $p=0.9$), renormalise, sample.  The set size adapts
dynamically: large when the distribution is flat, small when peaked.
}

% ==============================================================================
\chapter{Retrieval-Augmented Generation}
% ==============================================================================

\section{Motivation}

LLM parameters encode knowledge frozen at training time and can hallucinate.
\textbf{RAG} (Lewis et al.\ 2020) augments the LLM with a dynamically
retrieved, up-to-date knowledge base.

\defn{RAG Pipeline}{
\begin{enumerate}
  \item \textbf{Index}: encode all documents with a dense encoder;
        store embeddings in a vector database (FAISS, HNSW).
  \item \textbf{Retrieve}: embed the query $q$; find top-$k$ passages by
        maximum inner product: $\mathrm{score}(q,d) = E_Q(q)^\top E_P(d)$.
  \item \textbf{Generate}: concatenate query and retrieved passages as context;
        generate the answer with the LLM.
\end{enumerate}
}

\defn{Dense Passage Retrieval (Karpukhin et al.\ 2020)}{
Train dual encoders $E_Q, E_P$ with in-batch negatives:
\[
    L = -\log\frac{\exp(E_Q(q)^\top E_P(d^+))}
    {\exp(E_Q(q)^\top E_P(d^+)) + \sum_j\exp(E_Q(q)^\top E_P(d_j^-))}.
\]
}

% ==============================================================================
\chapter{LLM Agents}
% ==============================================================================

\section{Agent Architecture}

An \textbf{LLM agent} uses a language model as a reasoning engine that
perceives observations, maintains memory, plans, and executes actions via
external tools (web search, code interpreter, databases, APIs).

\defn{ReAct Framework (Yao et al.\ 2023)}{
Interleave \emph{Reasoning} (Thought) and \emph{Acting} (Action) steps:
\begin{enumerate}
  \item \textbf{Thought}: the LLM generates a natural-language reasoning step.
  \item \textbf{Action}: call a tool (search, calculator, code runner).
  \item \textbf{Observation}: the tool returns a result.
  \item Repeat until a final answer is produced.
\end{enumerate}
The reasoning trace makes the agent's decision process interpretable and
debuggable.
}

\section{Memory Types}

\begin{itemize}
  \item \textbf{In-context}: past conversation and retrieved documents within
        the prompt window.
  \item \textbf{External}: vector database of past experiences retrieved by
        embedding similarity.
  \item \textbf{Parametric}: knowledge encoded in model weights at pretraining.
\end{itemize}

\section{Multi-Agent Systems}

A \emph{planner} agent decomposes a task into subtasks; specialised
\emph{worker} agents execute each; a \emph{critic} agent reviews and
flags errors.  Communication is mediated through shared message queues or
structured prompts (e.g.\ AutoGen, CrewAI frameworks).

% ==============================================================================
% PART VII: ML ENGINEERING
% ==============================================================================

\part{ML Engineering}

% ==============================================================================
\chapter{Transfer Learning and Fine-Tuning}
% ==============================================================================

\section{Transfer Learning}

\defn{Transfer Learning}{
Pre-train on a large source task; adapt to a target task with limited
labelled data.
\begin{itemize}
  \item \textbf{Feature extraction}: freeze pre-trained weights; train only
        a new task head.  Fast; best when source and target are similar.
  \item \textbf{Full fine-tuning}: continue training all weights on the
        target dataset with a small learning rate.  Better when sufficient
        target data is available.
\end{itemize}
}

\section{Parameter-Efficient Fine-Tuning}

Fine-tuning billion-parameter models is expensive.  PEFT methods adapt a
small parameter subset.

\defn{LoRA (Hu et al.\ 2022)}{
Constrain the weight update $\Delta W$ to a low-rank factorisation:
\[
    \Delta W = A B^\top,\quad A\in\R^{m\times r},\; B\in\R^{n\times r},
    \quad r \ll \min(m,n).
\]
Only $A$ and $B$ are trained ($W_0$ frozen), reducing trainable parameters
from $mn$ to $r(m+n)$.  LoRA adds negligible inference latency when
merged: $W = W_0 + AB^\top$.
}

\textbf{Prompt tuning} (Lester et al.\ 2021): prepend a small set of
learnable continuous ``soft prompt'' vectors to the input embeddings; freeze
all other parameters.  Competitive with full fine-tuning at large model size.

% ==============================================================================
\chapter{Hyperparameter Optimization}
% ==============================================================================

\section{Problem Statement}

\defn{HPO}{
Find $\lambda^* = \arg\min_{\lambda\in\Lambda} f(\lambda)$ where
$f(\lambda)$ is validation loss after training with hyperparameters $\lambda$
(learning rate, architecture width/depth, regularisation, batch size).
Each evaluation requires a complete training run — expensive black-box
optimisation.
}

\section{Search Strategies}

\textbf{Grid search}: exhaustive enumeration of a discrete grid.  Exponential
in the number of hyperparameters; impractical beyond 3--4 dimensions.

\textbf{Random search} (Bergstra \& Bengio 2012): sample $\lambda$ uniformly
at random.  More efficient when only a few hyperparameters are important.

\defn{Bayesian Optimisation}{
Build a probabilistic surrogate model of $f$ (Gaussian process or Tree
Parzen Estimator).  Select the next configuration by maximising an
\emph{acquisition function}, e.g.\ Expected Improvement:
\[
    \mathrm{EI}(\lambda) = \mathbb{E}[\max(f^*-f(\lambda),\,0)].
\]
After each evaluation, update the surrogate.  Achieves sample efficiency
superior to random search for moderate-dimensional, expensive objectives.
}

\section{Multi-Fidelity Methods}

\textbf{Successive Halving}: run $n$ configurations for a small budget;
keep the top half; double the budget; repeat until one remains.

\textbf{Hyperband}: run successive halving with several different initial
bracket sizes in parallel to balance exploration and exploitation.

% ==============================================================================
\chapter{Distributed Training}
% ==============================================================================

\section{Parallelism Strategies}

\textbf{Data parallelism}: replicate the model on each device; partition
the mini-batch across devices; synchronise gradients via all-reduce.
Preferred when the model fits on one device.

\textbf{Pipeline parallelism} (GPipe): split layers into stages across
devices; pass micro-batches through the pipeline to keep all stages busy.

\textbf{Tensor parallelism} (Megatron-LM): split individual weight matrices
across devices; each computes a partial matrix product and results are
summed by all-reduce.  Required for very large layers.

\section{Gradient Synchronisation}

\defn{Ring All-Reduce}{
In a ring of $N$ workers, each sends and receives $1/N$ of the gradient
tensor per step.  Total communication $2(N-1)/N\,|\theta|$.  Bandwidth-
optimal: the bottleneck does not grow with $N$.
}

\section{Mixed-Precision Training}

Store master weights in FP32 but compute forward/backward passes in
FP16 or BF16.  Use \emph{loss scaling} to prevent gradient underflow.
Halves memory and doubles throughput on Tensor Core hardware.

% ==============================================================================
\chapter{AutoML and Knowledge Distillation}
% ==============================================================================

\section{Neural Architecture Search}

\defn{NAS}{
Automatically search a discrete space of architectures (layer types, widths,
skip connections) to maximise validation performance.
\begin{itemize}
  \item \textbf{RL-based}: a controller RNN proposes architectures as
        actions; validation accuracy is the reward.
  \item \textbf{DARTS} (Liu et al.\ 2019): relax discrete choices to a
        softmax mixture; optimise mixing weights $\alpha$ by gradient
        descent on validation loss; discretise at the end.
\end{itemize}
}

\section{Knowledge Distillation}

\defn{Knowledge Distillation (Hinton et al.\ 2015)}{
Transfer knowledge from a large \emph{teacher} $T$ to a smaller
\emph{student} $S$ by matching soft output distributions:
\[
    L_{\mathrm{KD}} = (1-\alpha)\,L_{\mathrm{CE}}(y,\hat y_S)
    + \alpha\,\tau^2\,D_{\mathrm{KL}}\!\left(
        \mathrm{softmax}\!\tfrac{z_T}{\tau}\;\Big\|\;
        \mathrm{softmax}\!\tfrac{z_S}{\tau}
    \right),
\]
where temperature $\tau>1$ softens the teacher's distribution, revealing
inter-class similarity structure not present in hard labels.
}

% ==============================================================================
% PART VIII: REINFORCEMENT LEARNING
% ==============================================================================

\part{Reinforcement Learning}

% ==============================================================================
\chapter{Multi-Armed Bandits}
% ==============================================================================

\section{The Bandit Problem}

\defn{$K$-Armed Bandit}{
At each round $t$, an agent selects arm $A_t\in\{1,\ldots,K\}$ and receives
stochastic reward $R_t\sim\mathcal{D}_{A_t}$ with unknown mean $\mu_a$.
Cumulative regret over $T$ rounds:
\[
    \mathrm{Regret}(T) = T\mu^* - \mathbb{E}\!\left[\sum_{t=1}^T R_t\right],
    \quad \mu^* = \max_a\mu_a.
\]
}

\section{Exploration--Exploitation Strategies}

\textbf{$\varepsilon$-greedy}: with probability $\varepsilon$ pull a random
arm; otherwise pull the empirically best arm.  Simple but suboptimal.

\defn{UCB1 (Auer et al.\ 2002)}{
\[
    A_t = \arg\max_a\!\left(\hat\mu_a + \sqrt{\frac{2\ln t}{N_a}}\right),
\]
where $N_a$ is the pull count for arm $a$.  The confidence bonus drives
exploration of under-tried arms.  UCB1 achieves $O(\ln T)$ regret.
}

\defn{Thompson Sampling}{
Maintain posterior $P(\mu_a|\text{history})$.  Each round, sample
$\tilde\mu_a\sim P(\mu_a|\cdot)$ and pull $\arg\max_a\tilde\mu_a$.  With
Beta--Bernoulli conjugacy this achieves near-optimal expected regret and
is empirically robust.
}

% ==============================================================================
\chapter{Markov Decision Processes}
% ==============================================================================

\section{MDP Framework}

\defn{Markov Decision Process}{
A tuple $(\mathcal{S},\mathcal{A},P,R,\gamma)$: state space, action space,
transition $P(s'|s,a)$, expected reward $R(s,a)$, discount $\gamma\in[0,1)$.
The agent follows policy $\pi(a|s)$ maximising expected discounted return
$G_t=\sum_{k\geq0}\gamma^k R_{t+k+1}$.
}

\defn{Value Functions}{
\begin{align*}
    V^\pi(s) &= \mathbb{E}_\pi\!\left[\sum_{k\geq0}\gamma^k R_{t+k+1}\mid S_t=s\right],\\
    Q^\pi(s,a) &= \mathbb{E}_\pi\!\left[\sum_{k\geq0}\gamma^k R_{t+k+1}\mid S_t=s,A_t=a\right].
\end{align*}
}

\section{Bellman Equations}

\prop{Bellman Expectation Equations}{
\[
    V^\pi(s) = \sum_a\pi(a|s)\sum_{s'}P(s'|s,a)[R(s,a)+\gamma V^\pi(s')],
\]
\[
    Q^\pi(s,a) = \sum_{s'}P(s'|s,a)\!\left[R(s,a)
    +\gamma\sum_{a'}\pi(a'|s')Q^\pi(s',a')\right].
\]
}

\prop{Bellman Optimality Equations}{
\[
    V^*(s) = \max_a\sum_{s'}P(s'|s,a)[R(s,a)+\gamma V^*(s')],
\]
\[
    Q^*(s,a) = \sum_{s'}P(s'|s,a)[R(s,a)+\gamma\max_{a'}Q^*(s',a')].
\]
Optimal policy: $\pi^*(s)=\arg\max_a Q^*(s,a)$.
}

% ==============================================================================
\chapter{Dynamic Programming}
% ==============================================================================

\section{Policy Evaluation and Improvement}

\defn{Iterative Policy Evaluation}{
Initialise $V_0\equiv0$.  Iterate:
\[
    V_{k+1}(s) = \sum_a\pi(a|s)\sum_{s'}P(s'|s,a)[R(s,a)+\gamma V_k(s')].
\]
Converges to $V^\pi$ (contraction mapping in $\ell_\infty$).
}

\defn{Policy Improvement}{
The greedy policy $\pi'(s)=\arg\max_a Q^\pi(s,a)$ satisfies
$V^{\pi'}(s)\geq V^\pi(s)$ for all $s$ (\emph{Policy Improvement Theorem}).
}

\defn{Policy Iteration}{
Alternate policy evaluation and policy improvement until convergence.
Terminates in finitely many steps for finite MDPs.
}

\defn{Value Iteration}{
Directly apply the Bellman optimality operator:
\[
    V_{k+1}(s) = \max_a\sum_{s'}P(s'|s,a)[R(s,a)+\gamma V_k(s')].
\]
Converges to $V^*$ without a separate policy evaluation step.
}

% ==============================================================================
\chapter{Temporal Difference Learning}
% ==============================================================================

\section{TD Prediction}

\defn{TD(0)}{
Update the value estimate after each transition $(S_t,A_t,R_{t+1},S_{t+1})$:
\[
    V(S_t) \leftarrow V(S_t) + \alpha\bigl[
    \underbrace{R_{t+1}+\gamma V(S_{t+1})}_{\text{TD target}}
    - V(S_t)\bigr].
\]
Bootstraps from current estimates (unlike Monte Carlo which waits for
episode end); can learn online and in continuing tasks.
}

\section{Q-Learning and SARSA}

\defn{Q-Learning (Watkins 1989)}{
Off-policy TD control:
\[
    Q(S_t,A_t) \leftarrow Q(S_t,A_t) + \alpha\bigl[
    R_{t+1}+\gamma\max_{a'}Q(S_{t+1},a') - Q(S_t,A_t)\bigr].
\]
Converges to $Q^*$ regardless of the behaviour policy (off-policy).
}

\defn{SARSA}{
On-policy TD control using the \emph{actual} next action $A_{t+1}$:
\[
    Q(S_t,A_t) \leftarrow Q(S_t,A_t) + \alpha\bigl[
    R_{t+1}+\gamma Q(S_{t+1},A_{t+1}) - Q(S_t,A_t)\bigr].
\]
Converges to $Q^{\pi_\varepsilon}$ (the $\varepsilon$-greedy policy).
}

\section{Deep Q-Network}

\defn{DQN (Mnih et al.\ 2015)}{
Approximate $Q^*$ with a deep network $Q_\theta(s,a)$.  Two stabilisation
techniques:
\begin{itemize}
  \item \textbf{Experience replay}: store $(s,a,r,s')$ in a replay buffer;
        sample random mini-batches to break temporal correlations.
  \item \textbf{Target network}: a periodically-frozen copy $Q_{\theta^-}$
        provides stable TD targets:
        \[
            L(\theta)=\mathbb{E}\!\left[
            \bigl(R+\gamma\max_{a'}Q_{\theta^-}(S',a')
            -Q_\theta(S,A)\bigr)^2\right].
        \]
\end{itemize}
}

% ==============================================================================
\chapter{Policy Gradient Methods}
% ==============================================================================

\section{Policy Gradient Theorem}

\defn{REINFORCE (Williams 1992)}{
Parametrise $\pi_\theta(a|s)$ directly.  The policy gradient theorem:
\[
    \nabla_\theta J(\theta)
    = \mathbb{E}_{\tau\sim\pi_\theta}\!\left[
      \sum_{t=0}^T \nabla_\theta\log\pi_\theta(A_t|S_t)\cdot G_t
    \right].
\]
Unbiased but high variance; reduces variance by subtracting a baseline $b$:
$\hat A_t = G_t - b(S_t)$ (the \emph{advantage}).
}

\section{Actor-Critic}

\defn{Actor-Critic}{
Maintain two networks:
\begin{itemize}
  \item \textbf{Actor} $\pi_\theta$: updated by policy gradient using
        the TD advantage $\hat A_t = R_{t+1}+\gamma V_\phi(S_{t+1})-V_\phi(S_t)$.
  \item \textbf{Critic} $V_\phi$: updated by TD learning.
\end{itemize}
The critic's bootstrap estimate reduces variance vs.\ REINFORCE while
remaining lower bias than using a fully-learned value function.
}

\section{Proximal Policy Optimisation}

\defn{PPO (Schulman et al.\ 2017)}{
Clip the policy ratio to prevent large policy updates:
\[
    L^{\mathrm{CLIP}}(\theta) = \mathbb{E}\!\left[\min\!\left(
      r_t(\theta)\hat A_t,\;
      \mathrm{clip}(r_t(\theta),1-\varepsilon,1+\varepsilon)\hat A_t
    \right)\right],
\]
where $r_t(\theta)=\pi_\theta(A_t|S_t)/\pi_{\theta_\mathrm{old}}(A_t|S_t)$.
Clipping prevents the new policy from deviating too far from the old one.
PPO is the dominant RL algorithm used for RLHF alignment of LLMs.
}

\end{document}
